\documentclass[UTF8,fontset=fandol,11pt,a4paper]{ctexart}
\usepackage[margin=22mm]{geometry}
\usepackage{amsmath,amssymb,amsthm,booktabs,array,longtable,xurl,tikz}
\usepackage{pgfplots}
\pgfplotsset{compat=1.18}
\usepackage[hidelinks]{hyperref}
\usetikzlibrary{arrows.meta,positioning}
\newtheorem{proposition}{命题}
\newcommand{\norm}[1]{\left\lVert#1\right\rVert}
\newcolumntype{L}[1]{>{\raggedright\arraybackslash}p{#1}}
\setlength{\parskip}{3pt}
\setlength{\tabcolsep}{4pt}
\setlength{\emergencystretch}{2em}
\renewcommand{\arraystretch}{1.12}
\raggedbottom
\title{移动 EV 主动拓扑辨识与源项 WZZT\\[3pt]\large 数学证明、可执行设计与地点--时段配对试验}
\author{Topo 研究实现：当前打开文档的续修版}
\date{2026 年 10 月 11 日续修：共同源、历史联合拟合与局部严格算例}
\begin{document}
\maketitle
\begin{abstract}
将移动 EV 定义为受接入、能量与控制能力约束的主动激励资源，分别建模 WZZT 结构拟合、试验设计与置信认证，并提出连接方式。
历史四案例的376个错误交换下界全部平凡，这项负面结果继续保留。
本次改为逐候选整盒AC相容性LP：共享R/X跨时重拟合，计入P/Q误差、未知源和非线性余项；用观测局部电流界与参数相关仿射上界增强排除，使用精确有理数校正对偶下界。
四例宽阻抗盒与测量误差AC算例的同96点观测中，参数仿射局部外包络排除277/376个错误候选，全局包络排除0/376个。固定soumalas11背景机制的20独立种子中，9 kWh PQ probing相对无EV的平均错误候选排除率差为+54.6个百分点；策略差额只按实际匹配对照解释。覆盖结论仅在物理域、源预算和条件噪声假设内成立，并限定于参考树加末端交换候选库；时窗、恒充、源项与盲端口的零/负结果均保留，尚未实现完整拓扑域的闭环WZZT认证。
保留以前的固定模板EV试验作为另一适用域的证据，不以其高成功率替代当前连续参数域验证。
结论限于声明的有限拓扑银行、完整连续参数盒、源预算与量测条件；真候选保留保证不等于本次输出的后验正确概率。
EV结果来自平衡等值AC和声明设备域；普通单/三相接口的最小电流、独立Q能力及实际跟踪仍需分别验收。
10月11日另给出三末端线性机制的精确证明与204次小LP核验：在固定公共干线源的位置前提下，共同源和历史共享参数保留了P-only激励的局部收益，同9 kWh调制使最小错误拓扑距离加倍，而逐输出源盒全部未决。该局部正结果不替代v5整盒AC四案例未知源位置下的负结果。
\end{abstract}

\section{研究定位与应当证明的对象}
合理背景是\textbf{受约束主动系统辨识、树响应核的层级恢复、带偏差的诚实推断}。
EV扩大可施加输入的地点集合；WZZT用层状原子与源路径拟合未知结构；置信模块决定能够保留或排除什么；设计模块决定下一批数据怎样采集。
主动逆变器 probing 已有可观测性和MILP研究\cite{cavraro,taheri}，墙盒充电记录也已有拓扑估计研究\cite{hotz}。
本研究的价值应由源误差、可执行移动动作、有效选择与未决规则的结合证明，不以“EV用于拓扑辨识”这一概念本身声称新颖。

\begin{center}\small
\begin{tabular}{L{3.1cm}L{5.5cm}L{6cm}}\toprule
对象 & 本文可处理 & 不能据此推出\\\midrule
局部树支撑 & 末端响应列的严格区间分离 & 全部隐藏物理母线可恢复\\
模板置信集合 & 预声明候选参考分布中的真模板保留 & 未知连续阻抗全部拓扑保证\\
整盒观测LP & 枚举候选树的全部声明R/X及有界源解释 & 候选银行之外的未知树自动被覆盖\\
参数估计 & 被激励方向及有信息的组合参数 & 单点探测已精确恢复全网R/X\\
EV收益 & 相同服务预算下的配对成功率、区间和成本 & V2G或独立Q对所有场景必然更好\\
\bottomrule\end{tabular}
\end{center}
全文区分错误认证控制、最终停止以及停止后参数精度；三个结论需要不同条件。
本轮主线是：第2.1--2.4节解释旧证书的语义及平凡原因，第\ref{sec:observationlp}节给新的整盒LP，第\ref{sec:v5results}节给同域实验与反例。
最新的理论--算例闭环见第\ref{sec:localv6}节：明确源位置的可解析小模型、全参数重拟合证明、源盒和历史消融及负结果。它先检验机制，不声称选址择时算法或完整WZZT闭环已完成。
其后保留的固定模板试验各有自己的候选范围，不能与整盒实验混算成功率。

\subsection{已有研究给出的数学背景}
\begin{center}\small
\begin{tabular}{L{3.4cm}L{5.7cm}L{5.5cm}}\toprule
研究路线 & 可借用的结论/工具 & 本文要补足的条件\\\midrule
图算法与逆变器MILP\cite{cavraro,taheri} & 通过电压对注入的响应恢复树结构 & 可用末端覆盖、隐藏节点等价类、源误差与设备能力\\
逆变器输入设计\cite{bhela} & 在运行约束下选择有信息的注入 & 车辆SOC、启动、能量服务及可接入时窗\\
T型模型区分\cite{dette} & 试验目标可直接针对候选模型差异 & 未知阻抗竞争解释必须重新拟合；局部最小值不能证明分离\\
带偏差回归推断\cite{armstrong} & 偏差与随机项共同决定区间宽度 & 有效统一误差界、准确PQ及数值泄漏预算\\
PAC/共形与序贯推断\cite{wilks,vovk,universal,howard} & 参考域覆盖、校准条件及可选停止工具 & 整episode独立单位、行动改变状态及现场分布偏移\\
实车数据\cite{saxena,esser} & 验证响应延迟、充放电和Q(P)能力 & 具体车辆/接口的逐项验收，尚需本项目现场验证\\
\bottomrule\end{tabular}
\end{center}
最适合当前阶段的是：有限可执行码本解决设计搜索，凸对比优化减少偏差半宽，LP/MILP为竞争结构提供拟合与下界。覆盖证明依赖可行性与统计条件；优化最优性决定效率，不能单独产生置信度。

\section{置信度审查：语义、负面证据与处理结论}
\label{sec:confidenceaudit}
\textbf{核查结论：两项批评均应接受，但须限定对象。}
主线的bootstrap频率不能进入结构认证接口充当概率；四个合成AC算例的既有全域下界在其交换邻域内全部平凡。
后者否定的是该证书在这些数据上的实际排除能力，尚不否定后文的条件覆盖证明，也未证明数据本身不可辨识。
第2.1--2.4节保留10月8日对2026年10月3日既有AC审计的复核；当时没有重跑四案例或EV试验。
第2.5节起给出10月9日起的新算法、完整证明及新试验，不能混淆两个时间版本。

\subsection{三种语义与三套置信方法不是同一分类}
\begin{center}\small
\begin{tabular}{L{2.4cm}L{5.6cm}L{6.6cm}}\toprule
语义层 & 输出实际含义 & 允许用途与禁用解释\\\midrule
1. 算法稳定性 & 重采样中某支撑的出现频率$s_{\rm boot}$ & 搜索提示、优先级；不是正确概率，不能单独固定为认证支撑。\\
2. 模型兼容性 & 残差、可行拟合质量，或有有效零假设校准的兼容性$p$值 & 判断观测是否与某模型相容；高兼容性不排除另一拓扑，近似卡方参照也非精确$p$值。\\
3. 结构风险控制 & 在声明模型域与采样机制下，错误认证事件的概率上界 & 接受部分结构、保留候选或未决；不自动给出本次输出的后验正确概率。\\
\bottomrule\end{tabular}
\end{center}
“多套置信方法在相同数据上的代价”节中的旧版A（条件响应区间）、B（PAC包络）、C（动作最大值候选覆盖）是\emph{第三层的三种构造}，各自仍有不同条件，不能与上表三层一一对应。
第2.5节新增的G/L/A则只是同一观测LP的三种包络强度，共用一套覆盖证明；其中字母A不表示上述旧版置信方法A。
后验$\Pr(h\mid Y)$另需先验与似然；输出后的选择性风险$\Pr(\text{错}\mid\text{输出})$另需独立风险校准。
全文不把这两者从$s_{\rm boot}$、残差或$1-\alpha$直接换算出来。

主线\path{rnj_wzzt/pipeline.py}的\path{_select_boundary_blocks}把bootstrap支持出现频率写入\path{confidence}，并用于选择随后固定的初始支撑。
因此这个字段应读作\emph{稳定性分数}，目前不能驱动认证账本。
\path{estimation/ac_validation.py}的局部残差诊断属于第二层；本报告的独立研究模块不等于主线已完成第三层闭环。
生产流程仍需可撤销候选与独立证书接口，本次未改生产主线。

\textbf{本次接口修复：}\path{confidence_v4.py}保留旧\path{status}以兼容调用，新增明确的阈值构造状态、证据类型和保证范围。
\path{requirements_verified=False}、\path{guarantee_applicable=None}明确表示代码未核验参考分布、真候选存在或现场动作机制。
\path{assumptions}标为要求清单，并按split/PAC分别声明交换性或独立同分布要求；阈值计算完成不再被元数据包装成假设验收完成。
不输出拓扑后验正确率或选择后条件错误率。
旧\path{ac_certificate.py}的\path{certified}也仅表示条件判据通过；\path{adaptive.py}账本检查风险预算与状态，尚不核验真实域覆盖或完整数值泄漏修正。
因此本次元数据修复消除了易误读的表达，但没有宣称已完成生产级域验收。

\subsection{四案例为什么全部得到平凡界}
以$S_{\rm obs}\in\mathbb R^{2n\times N}$记原始P/Q输入，$b$为白化原始量测的联合噪声球半径。
这里的$r_{\min}=r_{\min,\Omega}/Z_b$是线路电阻下界的\emph{标幺值}，$Z_b=V_b^2/S_b$；不能把欧姆值直接代入。
主线灵敏度矩阵在外部带倍数2，与后文吸收倍数2的$r_C$不同。
该证书计算
\begin{align}
 \underline\gamma&=[\sigma_{\min}(S_{\rm obs})-e_{PQ}]_+,
 &e_{PQ}&=\sigma_{PQ,\max}b,\\
 d_{\rm low}&=\frac{[2\sqrt2\,r_{\min}\underline\gamma-C_{AC}]_+}{K},
 &B_{\rm loss}&=[d_{\rm low}-b]_+^2,\label{eq:auditbound}\\
 \mathrm{FitQ}_{\rm upper}&=(1+B_{\rm loss}/m)^{-1}.\nonumber
\end{align}
$C_{AC}\ge0$为声明参数域上的AC余项包络，$K>0$把允许P/Q与电压校正换算到量测范数；$m$为原始量测维数。
$\mathrm{FitQ}$只是拟合质量分数，\emph{不是无功Q，也不是拓扑概率}。

\paragraph{第一行：先检验输入是否提供所有方向的信息。}
$n$是已计量末端数，$N$是时刻数；$S_{\rm obs}$的前$n$行是各末端有功，后$n$行是无功，列对应时刻。
这里的全行最小奇异值定义为
\[
 \sigma_{\min}(S)=\inf_{\|u\|_2=1}\|S^Tu\|_2
 =\sqrt{\lambda_{\min}(SS^T)}.
\]
向量$u$表示一个P/Q组合方向。若存在非零方向在全部时刻都没有激励，上式为0；若$N<2n$，它必为0，不能误取经济型SVD中省略零方向后的最小返回值。
原始噪声向量$e$包括实际使用的P/Q/V及可选首端/总表通道，$W$按各通道标准差白化。
$\|We\|_2\le b$是一个\emph{共同原始噪声事件}，不是拟合后的残差阈值；独立、零均值、已知方差的高斯通道下可取$b=\sqrt{\chi^2_{m,1-\alpha}}$。
有相关性或系统偏差时必须重新给出有效联合事件，不能直接套这个卡方半径。

令$E_{PQ}=S_{\rm obs}-S_{\rm true}$，$\sigma_{PQ,\max}$为P/Q通道标准差上界。
共同噪声球给$\|E_{PQ}\|_F\le\sigma_{PQ,\max}b=e_{PQ}$，而谱范数不超过Frobenius范数。
对任意单位$u$用反三角不等式，再对$u$取下确界，得到
\[
 \sigma_{\min}(S_{\rm true})
 \ge[\sigma_{\min}(S_{\rm obs})-\|E_{PQ}\|_2]_+
 \ge\underline\gamma .
\]
符号$[x]_+=\max(x,0)$保证下界非负。
所以$\underline\gamma=0$只表示\emph{这次扣除噪声后没有认证到正的全方向激励}；并不等于真实输入没有信息，也不等于任何结构都无法区分。

\paragraph{第二行：从结构差异到可观测距离，再到损失下界。}
在线性近似$D=HS$中，$H=2[R\ X]$包含共享路径核。
对于该审计中具有相反四点结构的换接候选，正边电阻下界给结构差异尺度$2\sqrt2\,r_{\min}$；
结合$\|AS\|_F\ge\sigma_{\min}(S)\|A\|_F$，得到线性响应分离下界$2\sqrt2\,r_{\min}\underline\gamma$。
该系数属于被审计的结构判据，不是对任意两个拓扑无需条件即可套用的公式。
AC模型还包含非线性损耗，故须减去覆盖允许域的$C_{AC}$；若余项已大于线性分离，截断后仍为0。

令$\overline H$为相关候选灵敏度算子范数的共同上界，$\sigma_V$和$\sigma_0$分别为末端、首端电压标准差上界，
则该实现的转换常数为
\[
 K=\sqrt{(\overline H\sigma_{PQ,\max})^2
          +(2V_{\max}\sigma_V)^2
          +(2V_{\max}\sqrt n\,\sigma_0)^2}.
\]
电压平方映射在$[0,V_{\max}]$上满足$|v_1^2-v_2^2|\le2V_{\max}|v_1-v_2|$；
同一个首端误差会进入$n$个末端压降，所以出现$\sqrt n$。
精确首端取$\sigma_0=0$。除以$K$把平方电压分离换算成白化原始量测距离下界$d_{\rm low}$。
再减去观测相对真均值的噪声半径$b$，反三角不等式给最小白化残差至少$[d_{\rm low}-b]_+$；平方后才得到最优平方损失的下界$B_{\rm loss}$。
每次减法都在支付一种允许的不确定性，不能为得到正数随意略去。

\paragraph{第三行：单调变换并不产生概率。}
若错误模型的最优白化平方损失为$L^*$，已有$L^*\ge B_{\rm loss}$，而$f(L)=1/(1+L/m)$递减，故
$f(L^*)\le f(B_{\rm loss})=\mathrm{FitQ}_{\rm upper}$。
下界为0时，上界等于1，只是这个转换没有排除能力；它不表示错误树“有100\%置信度”。
局部优化得到的任一可行损失$L_{\rm feas}\ge L^*$方向相反，最多提供拟合质量的下界$f(L_{\rm feas})$，不能拿来代替拒绝所需的损失下界。

\begin{center}\small
\begin{tabular}{lrrrrr}\toprule
96点算例 & 交换数 & $\sigma_{\min}(S_{\rm obs})$ & $e_{PQ}$ & $\underline\gamma$ & 非平凡界数\\\midrule
paper15 & 93 & 0.0007515 & 0.0136666 & 0 & 0 \\
soumalas11 & 45 & 0.0007923 & 0.0118253 & 0 & 0 \\
flynn16 & 102 & 0.0006760 & 0.0140887 & 0 & 0 \\
pengwah18 & 136 & 0.0006197 & 0.0148966 & 0 & 0 \\
\bottomrule\end{tabular}
\end{center}
四例合计376个非等价错误交换，均为\path{excitation_not_certified}：第一步即得$\underline\gamma=0$，所以$B_{\rm loss}=0$、$\mathrm{FitQ}_{\rm upper}=1$。
无噪声输入虽数值满行秩，其最小奇异值仅约$3.64$--$6.49\times10^{-5}$，不能凭含噪矩阵满秩宣称激励已认证。
这个1表示\emph{未得到有信息的质量上界}，不表示错误图质量确为1。

进一步，paper15的32点、93个错误交换、188次同预算局部拟合中，最小错误可行损失从冻结R/X时的15432.501降至重拟合后的1831.906；正确树可行损失为492.919。
错误图重拟合均耗尽12次函数评估预算；最佳错误拟合质量仅被夹在约$[0.45085,1]$。
这些可行损失是最优损失的\emph{上界}，不能替代拒绝所需的全域下界。
即使参考图已被假定为真，现有证书仍未排除任何该邻域错误交换，更不能据此认证全部错误拓扑。

证据来自主线\path{outputs/ac_topology_quality/analytic_four_cases96.json}和\path{full_fit_audit/report.json}；本次保存原文件及SHA-256于研究目录\path{results_confidence_review_20261008/}。
这是项目合成AC参考算例，不是四个现场实测馈线；旧legacy文档中的阻抗放大倍率不据此移植到本次审计。

\subsection{为什么EV高成功率不能填补这个缺口}
四案例排除问题允许错误拓扑在整个声明R/X盒域内重拟合，并对P/Q计量误差计费；当前该证书不支持额外源项。
EV主试验则是十端口、三个固定阻抗模板、60个冻结动作的参考分布校准。
二者的参数域、误差处理与推断对象不同。固定模板的97\%或74.7\%唯一正确率不能证明四案例证书已经恢复功效。
覆盖定理可以成立而始终返回大集合；优化器可找到很小残差却仍不能排除错误结构。
应同时报告\emph{有效性、实际功效、适用域}，不以任何一个替代另外两个。

\begin{proposition}[候选覆盖到部分支撑的安全转换]\label{prop:commonclades}
在声明域内，若$\Pr(h_\star\notin\mathcal C(Y))\le\alpha$，令非空候选集的输出为
\[
 \widehat{\mathcal L}(Y)=\bigcap_{h\in\mathcal C(Y)}\mathcal L(h);
 \qquad\mathcal C(Y)=\varnothing\ \Longrightarrow\ \widehat{\mathcal L}(Y)=\varnothing.
\]
则$\Pr\{\widehat{\mathcal L}(Y)\nsubseteq\mathcal L(h_\star)\}\le\alpha$。
\end{proposition}
\begin{proof}
在$h_\star\in\mathcal C(Y)$事件上，交集必为$\mathcal L(h_\star)$的子集；空候选时不输出任何支撑。
所以输出任一错误支撑的事件包含于真候选被排除的事件，不需再为交集中的支撑数量分摊风险。
\end{proof}
本命题只继承输入候选集的原有域与概率保证。有限模板不因此变成完整未知参数模型域；bootstrap候选池也不自动满足前提。
空候选必须返回\path{model_mismatch_or_exclusion_event}并检查模型/数据，不能利用空集交的逻辑真空接受全部支撑。
这里$h_\star$为真实结构，$Y$为整批观测，$\mathcal C(Y)$为所有尚未排除的候选。
$\mathcal L(h)$表示树$h$中非平凡的后代末端集合，即各内部边对应的clade；$\widehat{\mathcal L}$只输出每一棵保留树都同意的部分。
事件$\widehat{\mathcal L}\nsubseteq\mathcal L(h_\star)$表示\emph{至少输出了一个错误支撑}，不是某个支撑单独出错。
这正是无需再按输出支撑数量Bonferroni的原因：真候选一旦保留，全部公共支撑同时正确。
可能出现保留很多树却没有任何公共非平凡支撑；这应输出结构未决，不能强选一个拟合最优树并继承集合的风险保证。

\subsection{从负面结果得到可检验的改进目标}
近期主线应是\textbf{带未决输出的局部支撑辨识}：用WZZT与bootstrap生成可撤销争议，按兼容性筛查数据，再由独立响应区间或已校准候选集给第三层结论。
有限模板的结论明确命名为\path{reference_domain}；对完整未知R/X与源域未完成排除时输出\path{unresolved}，而非现场拓扑概率。

对争议支撑$C$，可用事前声明的两侧完整允许量测域$\mathcal M_C(a),\mathcal M_{\neg C}(a)$定义
\[
 d_C(a)=\inf_{\mu\in\mathcal M_C(a),\,\nu\in\mathcal M_{\neg C}(a)}
 \|W_a(\mu-\nu)\|_2.
\]
$a$包括接入节点、时窗、实际P/Q曲线和量测安排；$C$是一个争议支撑。
$\mathcal M_C(a)$与$\mathcal M_{\neg C}(a)$分别包含有、无该支撑时，遍历\emph{全部声明参数、背景和源域}得到的无计量噪声响应均值。
$W_a$将不同单位/精度的输出放到统一的误差尺度。
下确界表示两种解释最接近时的距离；比较两个点估计的距离会漏掉更接近的竞争解释。
计量噪声在下面的$b_a$中另计，不应在均值域和噪声半径里未经说明地重复定义。
如果真均值属于$\mathcal M_C(a)$、噪声范数不超过$b_a$且有可靠$\underline d_C(a)>2b_a$，则反三角不等式给所有错误域残差至少$\underline d_C(a)-b_a>b_a$。
具体地，令$Y=\mu_\star+e$；任意错误域解释$\nu$都满足
\[
 \|W_a(Y-\nu)\|_2
 \ge\|W_a(\mu_\star-\nu)\|_2-\|W_ae\|_2
 \ge\underline d_C(a)-b_a.
\]
这是设计的\emph{充分可分条件}；实际认证应检查观测数据对完整错误域的残差下界是否超过有效阈值。
局部内层拟合只能给分离最小值的上界。若两域闭包相交，例如目标边可趋于零，则$d_C(a)=0$；改称“定向设计”不能自动消除平凡界。
反过来，两个均值域在某个别的状态相交，并不意味着\emph{这次}观测无法排除其中一个。
第2.5节的新LP正是改做观测相容性检验，没有声称已经求出一个处处为正的$d_C(a)$。

下一轮比较必须沿用同一复合参数域与明确的误差预算，依次记录探测前后$\underline\gamma$、AC余项、目标分离下界、非平凡认证比例及错误认证/未决率。
EV只能通过改善有用输入方向、移动到争议祖先路径上的接入点或避开高偏差时段来争取这些条件；不能预先承诺证书一定转为非平凡。
保持当前全局证书作为严格基线，同时研究更紧输入误差界及争议相关方向；用现场标定或独立历史置信域收缩参数域，不能为获得好结果事后删掉难解释。



\subsection{新的解法：直接检验观测与整盒参数是否相容}
\label{sec:observationlp}
\textbf{问题变为：某个错误树在允许的全部R/X、计量误差和源项内，是否还能解释这次观测？}
这属于有界误差下的集合成员辨识与模型否证\cite{cerone2014,scott2014}。
以下LP与证明是本研究的推导；没有把文献中的固定参数模型、MILP等价变换或OPF松弛精确性直接移植过来。
它不需要先证明整个P/Q矩阵的最小奇异值为正，但也不保证每个错误树都能被排除。
原因是\emph{排除某个结构}不等于\emph{恢复任意完整灵敏度矩阵}：一个受正边权与共享路径约束的候选，可能已在少数观测方向上自相矛盾。
例如15个末端的原始P/Q矩阵有30行，24个时刻不可能满行秩；这会阻断原全方向证书，却不妨碍新的逐候选相容性检验。

\paragraph{一个严格的秩亏而仍可辨树的例子。}
仅为说明必要条件的区别，先取无噪声、无源、忽略损耗的线性模型，设$n$个末端、$T$个时刻，$Q=\kappa P$、$\kappa\ge0$，
$P\in\mathbb R^{n\times T}$满行秩。$\mathcal P_i$为根到末端$i$的路径，$Z_b=V_b^2/S_b$为基准阻抗。
此时$[P^T,Q^T]^T$秩至多为$n<2n$，旧全行最小奇异值必为0。
但平方压降矩阵满足
\[
 D=2K_{\rm eff}P,\qquad
 (K_{\rm eff})_{ij}=\frac1{Z_b}\sum_{e\in\mathcal P_i\cap\mathcal P_j}(r_e+\kappa x_e),\qquad
 K_{\rm eff}=\tfrac12DP^T(PP^T)^{-1}.
\]
每边有效权重$(r_e+\kappa x_e)/Z_b>0$。
$d_{ij}=K_{ii}+K_{jj}-2K_{ij}$、$d_{0i}=K_{ii}$因此给出根与末端的加性树距离，此处$K=K_{\rm eff}$。
也可直接从核重建层级：对任一非根分叉$v$，取其两个不同子分支内的末端$i,j$，则$K_{ij}$恰为$v$的根路径深度；集合$\{k:K_{ik}\ge K_{ij}\}$正是$v$的全部观测后代。
遍历条目阈值即可得到所有非平凡clade；相邻深度差给有效边权。条件是所有叶均已观测，并已合并无观测分叉的串联边，即恢复观测缩减树。
物理根仍是已知标记；若根到首个分叉有共同干线，$n\ge2$时其总权重为$\min_{i\ne j}K_{ij}$，不能因其clade是全体末端就丢掉这段深度。$n=1$时只能恢复总根叶有效权重。
所以分别估计任意$R,X$所需的$2n$行激励，并不是这种结构恢复的必要条件；独立Q仍有助于分开电阻/电抗贡献及克服有限噪声。
这个例子不消除AC余项或计量误差；下面的外松弛专门负责把它们纳入保证。

\paragraph{符号、单位与范围。}
固定候选树$h$，根为$0$，已计量末端为$j=1,\ldots,n$，采样时刻为$t=1,\ldots,T$。
$\mathcal P_i$为根到末端$i$的边集合，$J_e$为边$e$下游已计量末端，$f\succ e$表示$f$严格位于$e$下游。
每边参数$\theta=(r_e,x_e)_{e\in E_h}$在整个采样批次保持不变；$r_e,x_e$用$\Omega$，
$Z_b=V_b^2/S_b$用$\Omega$，功率、幅值电压和电流用相应标幺值。
默认平衡、径向、串联支路、无对地并联导纳（无shunt）；全部非根叶节点已量测P/Q/V，隐藏节点注入由源预算覆盖。
工程域为$0<\underline r\le r_e\le\overline r$、$0\le\underline x\le x_e\le\overline x$、
$0<V_{\min}\le |V_k(t)|\le V_{\max}$，最后一项必须覆盖\emph{全部}节点及未知源接入点。
这不是只检查已量测末端就能自动验收的现场条件。

令$\mathcal E$是原始量测同时误差事件，$\Pr(\mathcal E)\ge1-\alpha$；其上
\[
 |\widetilde P_{jt}-P_{jt}|\le b^P_{jt},\quad
 |\widetilde Q_{jt}-Q_{jt}|\le b^Q_{jt},\quad
 |\widetilde V_{jt}-|V_j(t)||\le b^V_{jt},\quad
 |\widetilde V_{0t}-|V_0(t)||\le b^0_t .
\]
波浪号表示实际输入算法的观测，$b$是事先有覆盖依据的绝对误差半宽，不是拟合残差。
平方电压降$D_{it}=|V_0(t)|^2-|V_i(t)|^2$不是幅值电压差。
把$[\widetilde V-b^V,\widetilde V+b^V]$与$[V_{\min},V_{\max}]$相交，所得正区间端点记为$V^-_{it},V^+_{it}$；首端同理。
于是
\begin{equation}
 D^-_{it}=(V^-_{0t})^2-(V^+_{it})^2
 \le D_{it}\le
 D^+_{it}=(V^+_{0t})^2-(V^-_{it})^2 .\label{eq:v5dropbox}
\end{equation}
这一步把首端误差也计入。若交集为空，则物理域与该测量盒不相容；代码记录此状态，不用倒置区间制造数值证书。

由原始观测可构造数据局部功率界
\begin{equation}
 \overline S_{jt}=\min\left\{S^{\rm prior}_j,
 \sqrt{(|\widetilde P_{jt}|+b^P_{jt})^2+(|\widetilde Q_{jt}|+b^Q_{jt})^2}\right\}.
 \label{eq:v5localS}
\end{equation}
$S^{\rm prior}_j$为独立有效的工程上限。\emph{数据局部不等于事后偷看真值}：在同一个$\mathcal E$上，三角不等式逐样本保证$|P_{jt}+\mathrm iQ_{jt}|\le\overline S_{jt}$。
因此这里无需再分样本或分配风险；相反，用本次最优拟合附近的小R/X盒替代完整工程盒没有此包含证明。

\paragraph{源项与AC余项。}
源默认表示未知位置的未计量消耗，逐时总预算为$\sum_s P_{st}\le B_P(t)$、$\sum_s Q_{st}\le B_Q(t)$，分量非负。
允许源在原有隐藏节点或原边内部；边内接入只将原边串联分段，各段R/X非负且分量和等于原边R/X，不允许添加阻抗无界的新支路。
多个源满足$\sum_s|P_{st}+\mathrm iQ_{st}|\le B_P+B_Q$；单一集中源可改用$\sqrt{B_P^2+B_Q^2}$。
任意多源不能擅用后者：$(P_1,Q_1)=(1,0)$、$(P_2,Q_2)=(0,1)$的模长和为$2>\sqrt2$。
记有效的总模长预算为$B_S(t)$。对原边$e$内每一实际支段$a$（未细分时$a=e$），KCL和三角不等式给
\begin{equation}
 \ell_a(t)=|I_a(t)|^2\le\overline\ell_e(t)
 =\frac{\big(\sum_{j\in J_e}\overline S_{jt}+B_S(t)\big)^2}{V_{\min}^2}.
 \label{eq:v5current}
\end{equation}
这里直接界节点电流相量的和，不在负荷功率之外再次重复叠加线路损耗。

定义观测P/Q给出的线性压降及其输入误差半宽：
\begin{align}
 \phi_{it}(\theta)&=\frac2{Z_b}\sum_{e\in\mathcal P_i}
 \left(r_e\sum_{j\in J_e}\widetilde P_{jt}+x_e\sum_{j\in J_e}\widetilde Q_{jt}\right),\label{eq:v5phi}\\
 \beta_{it}(\theta)&=\frac2{Z_b}\sum_{e\in\mathcal P_i}
 \left(r_e\sum_{j\in J_e}b^P_{jt}+x_e\sum_{j\in J_e}b^Q_{jt}\right),\label{eq:v5eiv}\\
 s_{it}(\theta)&=\frac2{Z_b}\sum_{e\in\mathcal P_i}(r_eB_P(t)+x_eB_Q(t)).\label{eq:v5source}
\end{align}
$\phi,\beta,s$均为未知$\theta$的线性函数，单位为平方标幺电压。
$\beta$保留参数依赖，避免把所有误差先乘整盒最大阻抗；真实源的线性贡献$g_{it}$满足$0\le g_{it}\le s_{it}(\theta)$。
源位置在不同输出间的耦合被放松，这使排除更保守；LP可行不等价于真的存在某一个源位置解释全部数据。

先在实际物理树上写恒等式；存在边内源时，它是原候选树的串联细分树，支段记为$a$，路径记为$\mathcal P_i^{\rm phys}$。
无边内源时，两棵树相同。支路流恒等式\cite{farivar}为
$v_{\mathrm{par}(e)}-v_{\mathrm{child}(e)}=2(r_eP_e+x_eQ_e)/Z_b-(r_e^2+x_e^2)\ell_e/Z_b^2$，
其中$v_k=|V_k|^2$，$P_e,Q_e$是送端支路功率，不能与末端负荷$P_j,Q_j$混淆。
代入支路功率等于下游注入与下游线路损耗之和，得到
\begin{align}
 D_{it}&=\phi_{it}(\theta)+\xi_{it}+g_{it}+\eta_{it},\qquad |\xi_{it}|\le\beta_{it}(\theta),\label{eq:v5identity}\\
 \eta_{it}&=\frac1{Z_b^2}\left[
 \sum_{a\in\mathcal P_i^{\rm phys}}(r_a^2+x_a^2)\ell_a(t)
 +2\sum_{a\in\mathcal P_i^{\rm phys}}\sum_{b\succ a}(r_ar_b+x_ax_b)\ell_b(t)\right]\ge0 .\label{eq:v5eta}
\end{align}
余项单边非负依赖非负R/X及上述串联模型，不要求每条支路有功始终正向。
对于原边内部的源，给所有子段共同的原边电流上界；同一原边的自项与段间交叉项合为
$\sum_k r_k^2+2\sum_{k<l}r_kr_l=(\sum_k r_k)^2$，X项相同。
不同原边的交叉项也能合并，得到在\emph{原候选树}上计算的二次上界
\begin{equation}
 \eta_{it}\le F_{it}(\theta):=\frac1{Z_b^2}\left[
 \sum_{e\in\mathcal P_i}(r_e^2+x_e^2)\overline\ell_e(t)
 +2\sum_{e\in\mathcal P_i}\sum_{f\succ e}(r_er_f+x_ex_f)\overline\ell_f(t)\right].\label{eq:v5etaouter}
\end{equation}
不要求各段R/X比相同。原边内电流可能随子段变化，所以这里是上界，不能把原边误当成具有一个处处相同的真实电流。

\paragraph{三个强度不同、域相同的LP方案。}
方案G把式\eqref{eq:v5current}中的$\overline S$换为全时刻工程上限，并把余项阻抗全部换成盒上界；
方案L使用数据局部$\overline S$，余项仍取盒最大值；方案A进一步保留余项对未知阻抗的仿射上界。
三者均在全部时刻重拟合\emph{同一整盒}R/X，没有固定成仿真真值。
对$u\in[l,u^+]$、$v\in[l,u^+]$有
\begin{align}
 u^2&\le(l+u^+)u-lu^+,\label{eq:v5secant}\\
 uv&\le u^+u+lv-lu^+,\qquad
 uv\le lu+u^+v-lu^+.\label{eq:v5mc}
\end{align}
第一式由$(u-l)(u^+-u)\ge0$展开得到；后两式分别来自$(u^+-v)(u-l)\ge0$和$(u^+-u)(v-l)\ge0$。
这是盒域上的割线与McCormick上平面。
把式\eqref{eq:v5etaouter}中的平方项用割线，交叉项分别统一采用式\eqref{eq:v5mc}的第一或第二上平面，得两个仿射函数$U^1_{it}(\theta),U^2_{it}(\theta)$。
由于所有电流权重非负，\emph{各自}都满足$\eta_{it}\le U^k_{it}(\theta)$；可同时施加两条约束，无需选一个拟合点收缩参数域。
一般逐项组合还能更紧，当前实现只保留两个整体平面以控制规模。

定义最小额外加性松弛$\tau$：
\begin{equation}
 \begin{aligned}
 \tau_h^*=\min_{\theta,\tau}\quad&\tau\quad\text{s.t. }\theta\in\Theta_h,\quad\tau\ge0,\\
 \phi_{it}(\theta)-\beta_{it}(\theta)&\le D^+_{it}+\tau,\\
 -\phi_{it}(\theta)-\beta_{it}(\theta)-s_{it}(\theta)-U^k_{it}(\theta)&\le-D^-_{it}+\tau,
 \quad \forall i,t,k .
 \end{aligned}\label{eq:v5lp}
\end{equation}
方案G/L只用一个常量$U$，方案A用两个仿射$U^k$，因此都是真正的LP。
若改用乘性膨胀$\kappa\beta(\theta)$则出现$\kappa\theta$乘积，不能仍称LP。
$\tau$的单位为平方标幺电压，表达还需额外放宽多少才能相容；它不是原AC最优损失、后验概率或$p$值。

在精确算术、相同观测和相同误差盒下，三方案满足$\tau_A^*\ge\tau_L^*\ge\tau_G^*$：
数据局部电流界不超过全域电流界；盒内每个平方割线及交叉上平面不超过该项的盒最大值，故A的约束蕴含L，L的约束蕴含G。
这是\emph{不会削弱理想排除能力}的结论，严格增益仍须实验验证。
不同求解器乘子给出的非最优有效下界未必逐次单调；比较数值结果时同时查看可行上界和差距。

\textbf{A方案改善的具体机制。} 盒最大余项允许错误候选一面用较小阻抗拟合线性压降，一面借用最大阻抗下的大余项吸收差异。
仿射上界把两者重新系在同一个$\theta$上。例如归一化$r\in[0,1]$时，$r^2\overline\ell\le r\overline\ell\le\overline\ell$；拟合若要求$r=0.1$，仿射包络只允许$0.1\overline\ell$，常量包络却仍允许$\overline\ell$。
这不是用拟合点缩域，因为约束对整个原盒都有效。跨时共享$\theta$又使错误树不能为每个EV功率点各换一套阻抗；P/Q轨迹通过增加相互制约的观测发挥作用。

\textbf{有符号源扩展。}
若源表示光伏、容性无功或有正负号的源项误差，预算必须改为
$\sum_s|P_{st}|\le B_P$、$\sum_s|Q_{st}|\le B_Q$。
此时$-s_{it}(\theta)\le g_{it}\le s_{it}(\theta)$，式\eqref{eq:v5lp}第一条压降约束也要减去$s_{it}(\theta)$。
第二条不变，电流仍用$B_P+B_Q$（单源可用勾股界）。当前API用\path{source_sign_mode}显式选择，默认非负；不能把未计量发电塞进非负源模型。
该预算界住的是模型中尚未表示的\emph{全部实际源}。若有非零名义源而只给其误差上界，须先显式加入名义源响应，或把名义幅值与误差一起纳入预算。

\begin{proposition}[整盒AC模型的外包含与条件排除]\label{prop:v5outer}
在上述物理域、源预算及误差事件$\mathcal E$上，真候选$h_\star$的真实参数使式\eqref{eq:v5lp}以$\tau=0$可行。
因此只有经过数值验证的下界$\underline\tau_h>0$才排除$h$；$\underline\tau_h=0$、求解失败或超时无正有效下界均保留为未决。
\end{proposition}
\begin{proof}
式\eqref{eq:v5dropbox}含真压降，式\eqref{eq:v5localS}--\eqref{eq:v5current}含真注入与电流，
式\eqref{eq:v5eiv}含输入误差。
由式\eqref{eq:v5identity}及$0\le\eta\le U^k$、$0\le g\le s$，有
$\phi-\beta\le D\le\phi+\beta+s+U^k$，与$D^-\le D\le D^+$合并即得两行约束。
有符号源将左侧换为$\phi-\beta-s$，证明相同。
真参数属于整盒且跨时不变，故$\tau=0$同时满足所有约束，最优值为0，任何有效正下界都不可能排除真候选。
\end{proof}
反向不成立：LP丢掉了多个输出共用输入误差、源位置、精确支路电流等耦合，可行只表示\emph{本证书尚不能排除}。
两树在已探测端口上有相同响应，或参数/源域能吸收差异时，增加24点或提高充电功率仍可能无法唯一辨识。

\subsection{概率保证和LP对偶：每一步分别保证什么}
\paragraph{共同量测事件到候选覆盖。}
本轮仿真假定各原始随机误差边缘为零均值高斯，已知标准差上界$\sigma_k$，另有确定偏差界$a_k$。
对本批$M$个非零随机误差通道，选择$\kappa>0$使
\begin{equation}
 M\frac{\exp(-\kappa^2/2)}\kappa\le\alpha,
 \qquad b_k=a_k+\kappa\sigma_k .\label{eq:v5mills}
\end{equation}
因为标准正态$Z$满足
$\Pr(|Z|>\kappa)\le\sqrt{2/\pi}\exp(-\kappa^2/2)/\kappa$
（对正半轴积分用$1\le z/\kappa$），并集界给$\Pr(\mathcal E)\ge1-\alpha$。
省略$\sqrt{2/\pi}<1$使阈值略保守，避免常数的数值下舍入；误差通道不需要独立。
标准差和偏差界必须确实有效，不能用同批拟合残差估计后直接代入声称精确覆盖。
若只有硬仪表误差界，可直接输入而无需高斯假设；若方差上界本身以$1-\delta$校准，则总失效概率还要计入$\delta$。
工程参数域或源预算若也来自统计估计，其失效概率同样须进入联合风险预算；一个未经校准的WZZT源点估计不能被当成确定性源界。

\begin{proposition}[一次批次的错误排除与公共支撑风险]
设预声明候选银行$\mathcal H$包含真候选，所有候选参数域和源预算均按命题\ref{prop:v5outer}构造，
$\mathcal C(Y)=\{h:\underline\tau_h\le0\text{或未得可靠证书}\}$。
则$\Pr\{h_\star\notin\mathcal C(Y)\}\le\alpha$，非空集合的公共支撑继承同一错误风险上界。
\end{proposition}
\begin{proof}
在同一个$\mathcal E$上真参数给$\tau=0$可行解，所以$\{h_\star\notin\mathcal C\}\subseteq\mathcal E^c$。
公共支撑结论再应用命题\ref{prop:commonclades}。
无需按错误候选数量分摊$\alpha$，因为它们共用同一原始量测事件；这并不覆盖没有列入银行的真拓扑。
\end{proof}
这是频率意义的\emph{真候选保留保证}，不是$\Pr(h\mid Y)$，也不是输出唯一树之后的条件错误率。
本轮每个冻结动作批次取$\alpha=0.05$；四案例、多种策略、多次重复的全部输出并不因此获得整体95\%保证。
历史数据选动作、然后采新数据时，所需的是给定历史之后仍成立的噪声边缘模型。
若同一批数据用于挑选最好动作，须覆盖全部潜在动作或重新独立确认。
多轮实施可事先取$\alpha_s=\alpha_0/[s(s+1)]$，由$\sum_s\alpha_s=\alpha_0$控制任意轮次错误接受；没有足够分离时允许永不唯一停止。

\paragraph{可靠下界不是求解器状态。}
把LP写为$\min c^Tz$，$Az\le b$，$\ell\le z\le u$，其中$z=(\theta,\tau)$。
给$\tau$设有限上限不改变最优值：先选盒中点$\theta^m$，逐行精确计算可行的$\tau^m$，再取$u_\tau>\tau^m$；至少已有一个可行点不高于此上限。
对任意$y\le0$，令$\rho=c-A^Ty$，则
\begin{equation}
 L(y)=y^Tb+\sum_j\min\{\rho_j\ell_j,\rho_ju_j\}
 \le \min_{Az\le b,\,\ell\le z\le u}c^Tz .\label{eq:v5dual}
\end{equation}
\begin{proof}
因$Az\le b$且$y\le0$，$y^TAz\ge y^Tb$；所以
$c^Tz=y^TAz+\rho^Tz\ge y^Tb+\sum_j\min\{\rho_j\ell_j,\rho_ju_j\}$。
这个结论不要求$y$最优，甚至不要求$\rho=0$。
\end{proof}
代码将求解器行乘子投影至$y\le0$，使用存储浮点数的精确有理值计算式\eqref{eq:v5dual}，最后向下舍入。
物理约束转成浮点矩阵时，逐行把$\sum_j|A^{\rm float}_{ij}-A^{\rm exact}_{ij}|\max(|\ell_j|,|u_j|)$加到右端并向上舍入，先保证浮点LP也是外松弛。
正有效下界才计为排除；同时修正原始可行解并给可行上界，若下界超过上界则报错。
此数值过程验证\emph{声明输入和LP}，不会验证现场阻抗域、计量精度、负荷符号或车辆能力。

LP对偶变量在这里有明确用途：构成可核验的排除下界、指出当前哪些压降约束限制相容性。
在变量上下界固定、最优基在邻域保持不变时，参数灵敏度形式为
$\mathrm d\tau^*=y^{*T}(\mathrm db-\mathrm dA\,z^*)$（目标固定）；退化时乘子可不唯一，不能直接把其大小当节点概率或全局派车收益。
MILP的连续松弛对偶解释的是松弛；固定整数拓扑后的LP对偶解释的是该拓扑条件下的连续问题。
二者均不能自动成为原MILP的全局边际价格。
只有相容性目标与阈值对应、且已证明外包含关系时，\emph{松弛下界}才能用于原物理问题的安全排除；任意WZZT目标的正LP下界并不足够。
无需为了得到有意义的排除下界而先造一个新MILP。

\subsection{同域数值验证：四案例、匹配动作和反例}
\label{sec:v5results}
\textbf{量测精度须按实际单位理解。} 本轮基准随机电压标准差$10^{-4}$ pu，按400 V线电压基准相当于0.04 V；有功标准差$2\times10^{-4}$ pu在200 kVA基准下相当于40 W。
这些是已声明的合成噪声条件，未由现场仪表校准得到；系统偏差、时间平均及实际精度不同须重新给出有效误差盒。
下列0/1/2倍敏感性只覆盖所列变化，不能当成一般现场计量鲁棒性的证明。
表中小数和统计区间端点为近似显示；排除决策使用归档的完整数值证书，不使用表格舍入值。
\paragraph{完成的实验范围与条件。}

四例均完成：正式比较176条策略记录、桥接12条、时窗32条、源项16条；另完成固定背景机制的20种子80条重复。正式比较中每例含两个生成树、两个预定测试种子；训练种子101--106与测试分离。

宽阻抗盒为每边$r\in[0.001,0.3]\,\Omega$、$x\in[0,0.3]\,\Omega$，全节点$V\in[0.9,1.1]$ pu，末端$|S|\le0.1$ pu。每条边参数在全部时刻共享并允许在整个盒内重拟合。$P,Q,V$合成测量噪声标准差依次为$(2,1,1)\times10^{-4}$ pu；使用每次试验95\%同时坐标盒。覆盖保证仍以物理域、源预算和条件噪声模型成立为前提，不能将整批试验共同称为95\%覆盖。

候选域是资产参考树与所有互不等价的末端两线交换树；不覆盖所有径向树。每个错误候选的排除必须来自正的精确校正对偶下界，未决及求解失败均保留。这里未运行闭环WZZT重建，也没有获得全拓扑域恢复证明。

\paragraph{同观测的旧界桥接。}

G为全局功率包络，L为观测局部功率包络，A再加入参数仿射AC余项上界；三者逐字节共用同一96点观测和新坐标误差盒。旧高斯球与新坐标盒是不同噪声事件；桥接复用四例生成机制和交换邻域，未声称重放历史原始数据。表内为排除候选数/错误候选数。

\begin{center}\small
\begin{tabular}{lrrrrrr}
\toprule
算例 & 错误候选 & 旧$\gamma$ & G & L & A & 真树保留 \\
\midrule
soumalas11 & 45 & 0 & 0/45 & 27/45 & 38/45 & 3/3 \\
paper15 & 93 & 0 & 0/93 & 24/93 & 79/93 & 3/3 \\
flynn16 & 102 & 0 & 0/102 & 0/102 & 39/102 & 3/3 \\
pengwah18 & 136 & 0 & 0/136 & 41/136 & 121/136 & 3/3 \\
\bottomrule
\end{tabular}
\end{center}

$\gamma=0$说明旧激励链未给出正下界；A取得正排除界说明该旧链平凡不等于拓扑不可辨识。新界仍可能未决。

\paragraph{可执行EV策略与配对比较。}

除无EV外，各策略均为单车6 h、24槽、并网有功9 kWh；此处匹配并网电量，不保证相同最终SOC。恒充是能量匹配对照，无EV为零能量基线。实际执行轨迹检查变流器、爬坡、SOC与离场条件；潮流仅在槽平均执行PQ处求解，未模拟槽内动态电压。移动策略为2 h首端口、1 h零注入转移、3 h次端口，另留转移电量预算，路线可达性为声明假设。所选端口与时窗仅依赖候选几何及独立历史PQ。

11个策略标签实际为8种不同注入；三组相同试验已缓存，不能当作独立重复。下表为每例四条预定生成树/种子记录上的平均错误候选排除率，分母始终包括所有错误候选。

\begin{center}\small
\begin{tabular}{lrrrr}
\toprule
策略 & soumalas11 & paper15 & flynn16 & pengwah18 \\
\midrule
无EV & 44.4\% & 7.3\% & 0.0\% & 21.7\% \\
固定端口恒充 & 44.4\% & 8.3\% & 0.0\% & 19.1\% \\
固定端口P & 46.1\% & 12.9\% & 0.5\% & 22.1\% \\
所选端口P & 49.4\% & 12.1\% & 0.7\% & 30.9\% \\
所选端口PQ & 83.3\% & 49.5\% & 14.7\% & 57.0\% \\
单车移动P & 73.3\% & 22.3\% & 1.5\% & 39.2\% \\
所选端口恒充 & 23.9\% & 8.1\% & 0.0\% & 23.5\% \\
PQ曲线去Q & 47.8\% & 11.6\% & 0.2\% & 30.1\% \\
\bottomrule
\end{tabular}
\end{center}

\begin{center}\small
\begin{tabular}{lrrrr}
\toprule
配对差（百分点） & soumalas11 & paper15 & flynn16 & pengwah18 \\
\midrule
PQ减所选恒充 & +59.4 & +41.4 & +14.7 & +33.5 \\
仅改变端口 & +3.3 & -0.8 & +0.2 & +8.8 \\
同P轨迹加入Q & +35.6 & +37.9 & +14.5 & +26.8 \\
移动P减单端口P & +23.9 & +10.2 & +0.7 & +8.3 \\
\bottomrule
\end{tabular}
\end{center}

所选恒充相对无EV的差依次为-20.6、+0.8、+0.0、+1.8个百分点。仅改变端口的对照是p\_site\_only减p\_fixed，同P轨迹的加Q对照是pq\_design减pq\_p\_only。移动P与单端口P比较同时改变峰值、时序和1 h停顿，属于整套策略净效果，不能将差额归因于双地点。移动或恒充出现零/负收益时照实保留，不能概括为EV普遍提高辨识率。四例小网格比较仅作描述，不给跨网络总体显著性结论。

\paragraph{预定机制下的独立重复。}

soumalas11固定背景场景0中，20个独立种子各自先平均两个生成树，再比较PQ与无EV。每种子差为两个整数排除数差之和除以90，符号判断不使用浮点容差。正/零/负差为20/0/0，平均差54.6个百分点，范围[34.4,80.0]；单侧符号检验$p=9.54\times10^{-7}$。符号检验仅用非零差，下面的CP则将零/负差均计为非正。“PQ有正收益”的概率单侧95\% Clopper--Pearson下限为0.861。该概率只针对这一冻结背景与合成高斯生成机制，既非现场可靠率，也非拓扑后验概率。

\paragraph{时窗与源项的负结果。}

相同所选端口P曲线在四个6 h时窗分别检验，表为两个生成树的平均排除率。历史规则所选时窗与事后最大值分开记录；这些补充试验不能用于事后改选并沿用原单次95\%保证。

\begin{center}\small
\begin{tabular}{lrrrrrr}
\toprule
算例 & 窗0 & 窗1 & 窗2 & 窗3 & 预选 & 事后最大 \\
\midrule
soumalas11 & 63.3\% & 83.3\% & 77.8\% & 74.4\% & 0 & 1 \\
paper15 & 6.5\% & 62.9\% & 48.4\% & 24.2\% & 0 & 1 \\
flynn16 & 0.0\% & 0.0\% & 17.6\% & 0.0\% & 0 & 2 \\
pengwah18 & 17.3\% & 60.7\% & 67.6\% & 21.3\% & 0 & 2 \\
\bottomrule
\end{tabular}
\end{center}

预选窗并非事后最优的算例为soumalas11、paper15、flynn16、pengwah18；低负荷平方代理忽略背景PQ自身携带的辨识信息，不能据此声称择时算法已验证有效或最优。

真实非负未计量负荷容量为$B_P=0.01$、$B_Q=0.0035$ pu；推断不获知其位置。正确容量模式覆盖源项，遗漏容量模式刻意令源预算为零。遗漏模式的高排除率即使出现也不能计为可靠优势；其源域前提已失效。源预算通过KCL与电压项共同放宽包络，可能降低辨识力。

\begin{center}\small
\begin{tabular}{lrrrr}
\toprule
算例 & 正确cap真树 & 正确cap排除 & 遗漏cap真树 & 遗漏cap排除 \\
\midrule
soumalas11 & 2/2 & 0.0\% & 2/2 & 83.3\% \\
paper15 & 2/2 & 0.0\% & 2/2 & 53.2\% \\
flynn16 & 2/2 & 0.0\% & 2/2 & 16.2\% \\
pengwah18 & 2/2 & 0.0\% & 2/2 & 58.5\% \\
\bottomrule
\end{tabular}
\end{center}

补充匹配敏感性仅在soumalas11的场景1、种子601进行，分别改变噪声尺度或阻抗盒。工程盒$r\le0.06\,\Omega,x\le0.05\,\Omega$是独立声明的示例短线先验，未获现场资产核验；这些探索结果不参与显著性选择。

\begin{center}\small
\begin{tabular}{lrrrr}
\toprule
参数盒 & 噪声倍率 & 无EV排除 & PQ排除 & 真树保留 \\
\midrule
宽盒 & 0 & 96.7\% & 100.0\% & 4/4 \\
宽盒 & 1 & 0.0\% & 83.3\% & 4/4 \\
宽盒 & 2 & 0.0\% & 40.0\% & 4/4 \\
工程盒 & 1 & 83.3\% & 90.0\% & 4/4 \\
\bottomrule
\end{tabular}
\end{center}

\paragraph{保留集能支持的结构。}

取所有未排除候选clade的交集，去掉单末端和全末端平凡集合；空保留集不认证结构。以下为每例四次正式试验的公共非平凡clade数量范围，而非对某个最优拟合树的可信度。结论仅限当前候选库及正确域事件。

\begin{center}\small
\begin{tabular}{lrrr}
\toprule
算例 & 无EV公共clade & PQ公共clade & 两组错误clade总数 \\
\midrule
soumalas11 & 0--0 & 1--4 & 0 \\
paper15 & 0--0 & 0--1 & 0 \\
flynn16 & 0--0 & 0--0 & 0 \\
pengwah18 & 0--0 & 1--1 & 0 \\
\bottomrule
\end{tabular}
\end{center}

全部176条正式记录的错误公共clade总数为0，空保留集0条；未将公共交集转换为全拓扑唯一恢复结论。

\paragraph{两个最小反例。}

一条正电阻边内存在$-8$ kW有功源，而末端计量仍为$+1$ kW、$Q=0$，实际$V_1=1.000224921>V_0=1$。在零测量误差与整阻抗盒内，遗漏/非负源模式的slack下界均为$0.000512393$ pu$^2$，有符号预算模式上下界均为0，真实参数零slack逐行精确可行。前者排除违反了其无源/非负源前提；将“源误差”只解释成正负荷会错误覆盖发电或有符号误差。

另一例中端口3零注入、4/5完全对称；在无关端口1施加6 h/24槽、0/3 kW理想ABBA共9 kWh，两个争议树仍生成相同电压，两者LP均成功且下界为0。该静态反例对端口1任意功率序列成立，未增加设备动态仿真；零收益来自结构缺信息。

\paragraph{完整审计与结论边界。}

全部归档332条策略记录均进入审计；物理工程域及源预算成立324条，本次实际误差落在误差盒内328条，三者同时成立320条。逐项未通过计数为测量误差盒:4、源预算:8。不同LP问题批次中候选LP失败0次，真树被排除0条策略记录，其中三者同时成立的记录0条。正式比较真树保留176/176、正确唯一2/176。误差越盒属于允许发生的随机事件；未删除源失配、误差越盒、物理域未覆盖或未决试验，也未将缓存别名当作独立样本。

本实现证明的是在声明域内的安全外逼近与候选排除，并验证局部/参数仿射包络和特定EV策略能改善有限候选集筛选；设备场地接入、在线源预算校准、完整候选生成与闭环WZZT更新仍需后续实现和现场验证。


\paragraph{三相 P 接口的独立探索。}
该补充不并入主试验：仅取 soumalas11、新预定 seed 981、两个真实候选，沿用冻结地点/时窗、宽阻抗域和源项为零。
稳态功率接口下限为 $4.2\,\mathrm{kW}$、上限为 $6\,\mathrm{kW}$，比较无 EV、单个充电块与四组双槽脉冲；两活跃动作均在 6 h 内输入 9 kWh，未按辨识结果调整波形。
“双槽脉冲”表示连续两个ON命令槽；执行过程含爬坡，OFF命令后的槽也可能因降坡而有正平均P，证书使用实际槽均值而非理想矩形命令。
\begin{center}\small\begin{tabular}{llll}\hline
真实候选 & 动作 & 排除错误候选（共45个） & 保留真图\\\hline
0 & 无 EV & 2 & 是\\
0 & 单充电块 & 24 & 是\\
0 & 四组双槽脉冲 & 29 & 是\\
1 & 无 EV & 2 & 是\\
1 & 单充电块 & 23 & 是\\
1 & 四组双槽脉冲 & 28 & 是\\
\hline\end{tabular}\end{center}
\begin{center}\small\begin{tabular}{lrrrr}\hline
动作 & 命令峰/kW & 轨迹峰/kW & 槽均值峰/kW & 积分能量/kWh\\\hline
单充电块 & 5.1429 & 5.1429 & 5.1429 & 9.0000\\
四组双槽脉冲 & 4.5000 & 4.5000 & 4.5000 & 9.0000\\
\hline\end{tabular}\end{center}
两活跃轨迹均通过该接口模型的启停、功率、SOC及能量审计；最大爬坡率为 1.0 kW/min，SOC 范围为 $[0.3408,0.4833]$，终值均为 0.4833。无 EV 仅为观测基线。
六个算例的槽均值 AC 全母线电压范围为 $[1.005208,1.023181]$ p.u.；新增连接电压上界 $1.05$ p.u. 的槽检查越界为 0/6，未筛去任何记录。
在平衡、单位位移功率因数、稳态执行P跟踪命令及上述连接电压条件下，两条实际动作换算的每相电流下界分别为 7.070 A 与 6.186 A。
这两个数针对实际采用的5.142857/4.5 kW稳态水平；通用4.2 kW模型下限本身并不在1.05 p.u.电压上保证6 A。
$6\,\mathrm{A}$ 仅作为控制设定兼容性参考，不是任意瞬时实际电流的下限；启停坡段不作此声称。$1.05$ p.u. 条件只检查了模拟槽均值，不能推及整个 $[0.9,1.1]$ 域或连续时间。
这是最低电流相制问题的条件修复，不是实车验收，也不支持普通 EV 硬件可独立调 Q。两活跃曲线峰值不同，不能把差异只归因于调制频率；该新 seed 仅作描述性探索，不作显著性检验。


\subsection{文献对照、尚未解决的问题与后续算法}
本次补查集合成员辨识、主动故障诊断、仿射模型区分和逆变器设计的原文；3篇选定PDF归档于\path{papers/mobile_ev_probing/}，版本、SHA-256与已读章节见\path{manifest.json}。
Scott正式版读第2、3.2、4节（含定义1、定理3）；Ding的2019年v3作者稿读第II--IV-A节（含引理1、定理1）；Bhela的2019年v2作者稿读第II--III节。
Cerone读第II节及第III节框架，Farivar读Part I第III-A节式(13)--(18)，后二者为在线正文；不声称五篇全文逐页读完。
公开来源已满足当前证明核对，未使用HKU订阅登录。

\begin{center}\small
\begin{tabular}{L{3.1cm}L{5.6cm}L{5.9cm}}\toprule
研究背景 & 对当前工作的帮助 & 不能直接继承的结论\\\midrule
集合成员辨识\cite{cerone2014} & 误差界定义所有可行解释，再做外松弛 & 一般非凸min--max框架不等于本AC问题已有现成LP\\
Zonotope主动诊断\cite{scott2014} & 输出集合不相交可保证区分；统一分离失败后具体观测仍可能诊断 & 固定线性模式及扰动集合不涵盖未知R/X全域\\
仿射模型区分\cite{ding2019} & 区分“对全部扰动，总有某时刻分离”的量词顺序 & MILP等价需要原文矩阵/耦合条件；未知阻抗乘输入需要额外松弛\\
逆变器动作库\cite{bhela} & 先筛运行可行性，再选信息较强动作 & 原文非计量负荷探测期近似不变不能直接用于6h现场\\
支路流模型\cite{farivar} & 给出准确平方电压/电流恒等式 & 本文没有借用OPF凸松弛精确性保证\\
\bottomrule\end{tabular}
\end{center}

当前值得继续的路线分为三档，按新增复杂度排列：
\begin{enumerate}
\item \textbf{已实现：整盒观测LP+有限可执行动作。}
研究模块已实现LP与候选交集；生产WZZT联动尚未接入。
在联动方案中，预声明工程银行$\mathcal H_{\rm eng}$，WZZT负责排序和提出可撤销争议，LP负责排除，未决集合交集给部分结构。
仅因WZZT或bootstrap未选中就删除某候选，会破坏银行含真模型的前提；除非另有覆盖证明，否则这些候选只能延后计算并保留为未决。
WZZT给出的点估计可用于动作排序，不能自动缩小认证使用的R/X或源域。
现有几何选点与低负荷择时只用于低成本排序；将其与剩余候选的真实排除效果配对核查，不能叫全域最优试验设计。
\item \textbf{尚未实现：保留共享误差/源的耦合外松弛。}
逐输出误差盒丢失共享$P/Q$误差、同一个源位置/时间轨迹的信息。
固定或枚举源落点，对$w=r_e\Delta P_{jt}$等乘积保留有限盒McCormick变量，可减少错误解释；仍须证明外包含，不能把松弛可行点当物理可行点。
\item \textbf{尚未实现：困难候选的覆盖式分支细化与鲁棒设计。}
把完整参数/源域分成覆盖它的子盒，只有所有子盒都有正有效下界才排除父候选。
EV外层可最大化最坏竞争分离的可靠下界；内层局部拟合仅给分离最小值上界，不能充当保证。
超时保留未决，并记录箱数、上下界差和运行代价。
\end{enumerate}
\paragraph{待实现：保留固定单源的跨节点、跨时耦合。}
下面是自行推导的候选LP收紧方案，尚未实现或测试，不能声称解决了pengwah的源预算平凡结果。
额外前提是整个批次恰有一个未建模实际源，位置固定、线路参数不变，且逐时
$0\le p_t\le B_P(t)$、$0\le q_t\le B_Q(t)$；本段不覆盖位置变化、多源或有符号源误差。
在固定候选树中，源的线性平方电压响应可写为
\[
 g_{it}=a_i^Pp_t+a_i^Qq_t,
 \qquad 0\le a_i^P\le\frac2{Z_b}\sum_{e\in\mathcal P_i}r_e,
 \quad 0\le a_i^Q\le\frac2{Z_b}\sum_{e\in\mathcal P_i}x_e.
\]
这里$a_i^P,a_i^Q$跨时共享，$p_t,q_t$跨节点共享；$a$是源线性项系数，不是完整AC雅可比。
真实$a$由根到源与根到端口的共同路径决定，边内源还含共同路径的部分段阻抗。
仅施加上述上界会放松同一源位置的几何耦合，但仍包含任意固定节点/边内位置。

\begin{samepage}
取\emph{常数}工程上界
$\overline a_i^P=2|\mathcal P_i|\overline r/Z_b$、
$\overline a_i^Q=2|\mathcal P_i|\overline x/Z_b$。
对$w^P_{it}=a_i^Pp_t$，施加四条McCormick外约束
\begin{align}
 w^P_{it}&\ge0, &
 w^P_{it}&\le\overline a_i^Pp_t,\\
 w^P_{it}&\le B_P(t)a_i^P, &
 w^P_{it}&\ge\overline a_i^Pp_t+B_P(t)a_i^P-\overline a_i^PB_P(t).
\end{align}
\end{samepage}
Q同理，令$g_{it}=w^P_{it}+w^Q_{it}$。即使$B_P(t),B_Q(t)$随时刻改变，只要它们为外给常数，仍是LP。
不能把$\overline a_i$换成含未知$\theta$的路径和后直接乘未知$p_t$，否则重新出现双线性项。
原来的AC电流及余项包络继续使用整个单源预算，不能用本次拟合$p_t,q_t$缩小认证域。
两侧压降约束改为
\[
 \phi_{it}+g_{it}-\beta_{it}\le D^+_{it}+\tau,
 \qquad
 -\phi_{it}-g_{it}-\beta_{it}-U^k_{it}\le-D^-_{it}+\tau.
\]

\textbf{包含与强度。}
真实双线性点满足四条约束，所以固定单源真模型仍以$\tau=0$可行。
又有$0\le g_{it}\le B_P(t)a_i^P+B_Q(t)a_i^Q\le s_{it}(\theta)$，
故新LP每个可行点投影到$(\theta,\tau)$后必满足旧逐输出源包络；相同输入和参数域下，新最优slack不会小于旧值。
但严格增益不保证。考虑固定参数的纯P代数例子，归一化后$\overline a_i=B_t=1$且路径响应上界也为1，零输入误差时，若所需源残差为
$G=\bigl[\begin{smallmatrix}1&0\\0&1\end{smallmatrix}\bigr]$且$U=0$，旧盒允许它，
新约束则由两个对角1迫使$a_1=a_2=p_1=p_2=1$，再迫使非对角项至少为1，产生矛盾。
同一反例在余项统一上界$U<1/2$时仍矛盾。
这是松弛强度的代数例子，不是新增AC算例或性能证据。

\textbf{McCormick仍可能完全无效。}
同样归一化、纯P且$U=0$，取
$G=\bigl[\begin{smallmatrix}1/2&0\\0&1/2\end{smallmatrix}\bigr]$。
真实固定单源乘积$a p^\top$不能产生它（矩阵秩为2），
但$a_1=a_2=p_1=p_2=1/2$时，每个McCormick变量都可独立落在$[0,1/2]$，所以新LP仍有零slack。
更一般地，$a_i=\overline a_i/2$、$p_t=B_t/2$使整个半幅盒
$0\le w_{it}\le\overline a_iB_t/2$都通过松弛。
因此多个竞争树若都能将残差留在这一盒内，所有排除下界仍可能为0；不能将“共享变量”写成已恢复精确秩一耦合。

\textbf{未来实现步骤。}
\begin{enumerate}
\item 事先声明固定单源假设、逐时预算及全节点电压域；仍保留整盒R/X。
\item 建立跨时$a$、跨节点$p,q$及McCormick变量，保留原AC包络，求最小非负额外slack。
\item 用外向数值保护及精确校正对偶下界排除；下界0、失败或超时均未决。新的有限slack上限须从新LP的已验证可行点重新构造。
\item 可另枚举源节点/边段并加入路径响应等式；只有所有允许位置域均获正下界才排除候选。枚举节点不能声称覆盖未枚举边内位置。
\end{enumerate}
固定位置若有有效依据，收紧继承同一量测事件，无需另分一次风险；未经核实的固定单源假设则不能获得现场无条件保证。
有符号源可另用正负分量及绝对预算构造外松弛，多源需另作源编号与总预算耦合；不能直接套用本段模型。


若现场已有首端总P/Q表，应优先保留这些观测。本轮v5只使用首端电压；在当前无shunt模型中，总表有功减去末端有功之和等于全部未计量源有功加线路有功损耗。
因此总分表的同时误差界、损耗非负性及全域损耗上界，可给源幅值新的外区间，再与WZZT源拟合衔接；不能直接把拟合源均值当预算。
额外计量通道必须纳入原始联合事件，这项观测增强尚未在本轮测算，也不保证一定解决源包络平凡问题。

另一个可直接沿当前数学框架推进的方向是\textbf{参数/响应的投影区间}，尚未实现于本轮接口。
记$\mathcal F_h(Y)$为候选$h$在$\tau=0$时的LP参数可行集合，$m_h(\theta)$为所关注的线性响应量，例如共享路径核中某探测列的线性化电阻灵敏度（不是完整AC雅可比）。
对每个保留候选，分别求$\inf_{\theta\in\mathcal F_h}m_h(\theta)$和$\inf_{\theta\in\mathcal F_h}[-m_h(\theta)]$的可靠下界$L_h^+,L_h^-$，则
\[
 \left[\min_{h\in\mathcal C(Y)}L_h^+,\;
       \max_{h\in\mathcal C(Y)}(-L_h^-)\right]
\]
在同一个$\mathcal E$上包含真实响应，因真参数本就在某个保留候选的可行集中。
无需另作bootstrap，也不额外消耗一次量测事件的风险预算；但必须使用正确方向的优化\emph{下界}，不能把局部可行极小值误作区间下端。
未求解或求解失败的保留候选须使用整盒先验给出的线性量上下界，不能从区间并集中漏掉；空可行集也须经可靠证书确认后才能删除。
这样可将“排除错误拓扑”和“仍存结构下的参数精度”分开量化；尚未运行的投影区间不列为本轮结果。

\paragraph{择时反例之后，应优化什么？}
低负荷只可能减小式\eqref{eq:v5current}的余项，不能保证背景P/Q变化仍足以区分拓扑。
本轮全部末端背景P/Q已经量测并进入$\phi(\theta)$，本身就是可利用的辨识激励；未量测变化才由源预算或误差域承担。
因此，量测覆盖与拟合方式决定应如何权衡安静时窗和丰富输入，不能把所有背景变化都当成干扰。
后续动作$a$应同时指定接入点、时间窗和实际执行P/Q轨迹。设$\mathcal A_{\rm safe}$为满足设备、行驶及电网运行约束的有限动作库；当前试验尚未实现完整的电网运行安全筛选。
对每个保留候选，定义$\mathcal Y_h(a)$为该动作下全部允许参数、背景负荷和源项所产生的\emph{无计量噪声输出集合}，通道与未来相容性检验采用的原始P/Q/V及首端量测一致。
参数可取先前观测给出的外可行域$\mathcal F_h$，但未来背景及源的范围必须另有有效依据；有限历史样本本身不能保证包含未来状态。
固定一个事前声明的正定权矩阵$W_a$及范数$\|z\|_{W_a}=\sqrt{z^TW_az}$，令新批次噪声在其中的同时半径为$b_a$，可考虑
\begin{align}
 d_{hg}(a)&=\inf_{y\in\mathcal Y_h(a),\,z\in\mathcal Y_g(a)}\|y-z\|_{W_a},\label{eq:v5designsep}\\
 a^*&\in\arg\max_{a\in\mathcal A_{\rm safe}}
       \min_{h\ne g\in\mathcal C}\{\underline d_{hg}(a)-2b_a\},\qquad
       \underline d_{hg}(a)\le d_{hg}(a).\label{eq:v5designmaxmin}
\end{align}
式\eqref{eq:v5designmaxmin}用于$|\mathcal C|\ge2$；空集合报告模型不相容，单一候选已在声明银行内完成结构区分，不用空集合最小值制造“无限收益”。
这里$\underline d$必须是可靠下界。若最坏项为正，任意两候选的含噪输出集合不相交：
若同一观测同时距$y,z$不超过$b_a$，三角不等式将给$\|y-z\|_{W_a}\le2b_a$，与分离条件矛盾。
这是\emph{所有允许真实状态都能区分}的强目标，不是预期收益。
若$\mathcal Y_h$取精确物理输出域，分离只保证理想模型可区分，不保证较松的G/L/A检验也能排除错误树。
要保证具体LP的排除，须对该LP对应的输出外集合证明分离，或在新观测到达后取得实际的正对偶证书。
其中任何一对在该动作下结构性等价，最坏分离就可能为零；零值不表示该动作对本次观测必定无用。
这里零值指$d$，含噪裕量$d-2b_a$还要扣除噪声半径。
若历史外可行域的覆盖风险为$\alpha_{\rm hist}$，新批次事件风险为$\alpha_{\rm new}$，上述闭环保证至少要合计二者，不能在更新参数域之后只计算新噪声。

低成本替代是用\emph{预先冻结的训练场景}估计每个动作的排除比例，选定后再用新批次的观测LP认证。
前者只是分布相关的规划评分，后者才负责风险保证；可保留一次充电作为比较基线。
未知R/X与输入相乘、AC方程及不确定背景使式\eqref{eq:v5designsep}一般非凸；文献中的固定仿射模型MILP不能直接套用。
可先枚举有限动作，以参数覆盖分支和外松弛求可靠分离下界；局部拟合给出的距离\emph{上界}仅适于发现难区分竞争者。
以上是针对本轮择时反例提出的待实现算法，不是另一个已经验证优于基线的结果。

\textbf{单次无法统一分离，也不等于多轮无效。}
一个纯数学例子是：模型$h$在动作A/B下输出$(u,u)$，模型$g$输出$(v,1+v)$，且$u,v\in[0,1]$。
分别看A或B，两个输出区间都相交；联合两个动作并要求同一个参数跨轮共享，两线段的欧氏距离却为$1/\sqrt2>0$。
因此移动EV的潜在价值还包括从不同端口取得互补约束；这需要联合批次设计与累计风险，不能由单动作评分为零宣布永久不可辨识。

现阶段没有证据要求用一个大型MILP同时承担拓扑搜索、派车、源拟合和风险保证；分开后每一步更容易审计。
完整未知拓扑域、三相不平衡、配变与并联支路、现场计量误差标定、真实接入/路径约束仍未解决。
真实15min平均量测通常不满足$\mathrm{AC}(\overline P,\overline Q)=\overline{\mathrm{AC}(P,Q)}$。
现场应选择稳定驻留后的同步采样，或独立标定包含控制与平均化影响的误差预算；本轮静态AC槽平均注入试验不能替代该验证。

\paragraph{建议的闭环接口。}
一次流程依次为：声明工程候选/参数/源域与风险预算；WZZT排序争议但不未经认证地锁死支撑；
对原始观测运行整盒LP并保留未决；输出非空保留集合共有的支撑；
在仍有争议的接入端口上按设备、时窗、行驶及运行约束选择下一动作，采新数据后重复。
若所有单次合法动作均无正分离保证，应返回当前保证不足，并考虑联合动作或扩大观测覆盖；只有更强的联合等价证据才能宣布结构性不可辨识。
移动性提供改变探测端口的手段，独立Q提供第二种响应方向，源包络防止把遗漏负荷当结构差异，风险模块限制接受什么；四者各有作用，任何一项都不能替代其余前提。



\section{EV模型：车辆、接口与量测构成一个系统}
\textbf{本轮模型的相制边界。} 四例采用400 V、200 kVA平衡等值AC，EV总P/Q作为等值注入；本研究没有验证单相EV接入后的三相四线不平衡电压。
普通P策略采用1.4 kW最小非零功率，这接近单相230 V、6 A、单位功率因数时的1.38 kW。
制造商资料说明6 A充电电流设定下限及单/三相配置\cite{evminimum}；若实际三相设备每相最低设定6 A且稳态按设定跟踪，则400 V线电压下总有功约$\sqrt3\times400\times6/1000=4.16$ kW。
设定值不等于任意瞬间实测电流，启动、温度限制及接近充满时仍需实测跟踪或另行建模。
因此当前约3 kW的P波形不能直接宣称适用于这种普通三相充电接口。
必须按目标车/接口重设功率下限、相切换、启停和Q能力，或采用经过验证的可控变流接口；单相方案则须在不平衡模型中重新检验。
下文“可执行”均指通过所声明设备模型的合成检查，不能读成已经完成实车或逐相电网验收。
车辆须通过兼容并网接口接入已知节点/相别，并记录实际P/Q和同步电压。
DOE支持车辆与匹配EVSE作为移动储能平台\cite{doe}；这不意味着汽车到达任意末端就能接入。
普通AC EVSE提供交流与控制信号，变换器在车内，因而调节电流上限不等于可独立调Q\cite{evse}。

\begin{center}\small
\begin{tabular}{L{2.4cm}L{5.8cm}L{6.4cm}}\toprule
设备域 & 输入约束 & 所需硬件及辨识含义\\\midrule
单向P & $P\ge0,\ Q=\kappa P$ & 默认$\kappa=0$；一般估计$R+\kappa X$组合响应。应用正权树证书另需组合边权非负。\\
单向PQ & $P\ge\ell_Q(Q)$，Q独立 & 已验证无功能力的变换器；不从电池输出有功。\\
双向PQ & $-P_d\le P\le P_c$，Q独立 & 车辆、接口及通信同时支持双向和无功；V2G本身不证明独立Q。\\
\bottomrule\end{tabular}
\end{center}

正P表示从电网充电，$0<\eta_c,\eta_d\le1$。所有域满足
\begin{align}
 P^2+Q^2&\le S_{\max}^2,&|\dot P|&\le\rho_P,&|\dot Q|&\le\rho_Q,\\
 \dot E&=\eta_c[z]_+-[-z]_+/\eta_d,&z&=P-\ell_Q(Q),&\ell_Q(Q)&=k_QQ^2,\label{eq:evenergy}\\
 E(0)&=C_b\,\mathrm{SOC}_0-E_{\rm drive}-E_{\rm setup},&
 C_b s_{\min}&\le E(t)\le C_b s_{\max},&E(6h)&\ge C_b s_{\rm depart}.
\end{align}
历史v4试验设置电池容量$C_b=60$kWh、SOC初值0.3、离站至少0.4、域$[0.1,0.9]$、效率0.95、驾驶0.5kWh、换接辅助预算0.05kWh。
本轮v5改用SOC初值0.35、主P接口最小非零驻留1.4kW、响应延迟0；移动续段继承此前SOC，移动驾驶预算2kWh；三相补充另用4.2kW下限。各版本参数不能混用，完整配置见各自归档。
$S_{\max}=7.5$kVA，$P_c=6$kW、$P_d=3$kW、$|Q|\le4.5$kvar、两轴爬坡率1单位/min；$k_Q=0.002$kW/kvar$^2$是声明的研究损耗参数。
上述数值须用具体车辆/接口验收，尚非现场标定。

\subsection{连续轨迹与24个真实平均值}
每个15min槽由指令延迟、线性爬坡和驻留组成；末槽预留停机归零。
爬坡持续时间取
\[
\tau=\max\{|\Delta P|/\rho_P,|\Delta Q|/\rho_Q\}.
\]
若延迟、爬坡及停机不能放入槽长则拒绝。
\begin{proposition}[声明设备模型的连续段可行性]\label{prop:evsegment}
在对应设备模式中，若线性爬坡两端满足P/Q盒、功率圆及适用的$P\ge k_QQ^2$约束，且非零段持续时间$\Delta t\ge\tau$，则整段满足这些功率约束与爬坡率限制。
若还准确检查每段端点和电池功率$z$的全部零点处的能量界，则整段SOC可行。
\end{proposition}
\begin{proof}
盒、圆及凸函数$k_QQ^2$的上图均为凸集，合法端点的凸组合仍合法。爬坡斜率满足
\[
 |\dot P|=|\Delta P|/\Delta t\le\rho_P,\qquad
 |\dot Q|=|\Delta Q|/\Delta t\le\rho_Q,
\]
其中使用了$\Delta t\ge\tau$。
线性P/Q段上的$z=P-k_QQ^2$为二次多项式，$\dot E$与$z$同号，故内部SOC极值只能出现在$z=0$处，驻留段及端点另查。
\end{proof}
最小非零充电功率是非凸例外，逐段要求驻留为0或高于最小值，仅允许启动/停机爬坡穿越低功率区。
式\eqref{eq:evenergy}按分段多项式积分；时间须统一，分钟爬坡时长应除以60后用于kWh积分。实际审计使用浮点根和容差；指令延迟由轨迹生成器保证，尚无独立区间算术全过程安全证书。

历史v4主码本允许连续细调到近零P，这是可控接口假设。其普通AC车辆补充另设1.4kW最小驻留、1min延迟，生成0/3kW的ON--OFF控制；不能声称任意车辆允许频繁重启。
Saxena实测特定Leaf/充电器常规响应约秒级，但首次启动109.32s且极限指令饱和\cite{saxena}；Esser的8车142条记录揭示Q(P)及动态行为随车辆/接口变化\cite{esser}。
设备例证沿用前轮已读Saxena、Esser全文；Hotz墙盒研究仅按摘要/书目信息定位相关工作。本次重整为自行推导，未重新进行文献检索，也未把他人实车数据写成本项目实测。

\subsection{公平预算与额外硬件能力的成本}
全部活跃策略匹配6h、24点、9kWh网侧净充电及共同额定限制；实际峰值、吞吐、电池增能不同。
忽略Q损耗时，$\Delta E_b=\eta_c E_{\rm net}-(1/\eta_d-\eta_c)E_{\rm export}$；所以相同净电量的双向循环会减少车辆获得的能量。
以下为主研究冻结码的实际连续积分：
\begin{center}\small
\begin{tabular}{lrrrr}\toprule
控制 & 网侧净kWh & 电池增能kWh & 实际峰值kVA & 离站SOC\\\midrule
恒功率 & 9.00 & 8.55 & 1.51 & 0.433 \\
单向P & 9.00 & 8.55 & 2.97 & 0.433 \\
单向PQ & 9.00 & 8.36 & 5.37 & 0.430 \\
双向PQ & 9.00 & 7.69 & 7.37 & 0.419 \\
\bottomrule\end{tabular}
\end{center}
无车基线不交付9kWh服务，恒功率才是比较调制收益的匹配对照。恒功率仍包含真实启停，小幅边缘变化不能被说成严格零激励。
未知线路热额定值、变压器、三相及接入保护未进入本版事前网络安全约束；设备可执行不等于完成现场安全验收。

\clearpage
\section{数学主线、统一记号与假设}
以下命题是本文在明确假设下的自行推导，文献提供模型与方法背景；不将数值成功率写成定理。
两条证明链分别是
\[
\text{树响应模型}\ \Longrightarrow\ \text{同时响应区间}\ \Longrightarrow\ \text{局部支撑证书},
\]
\[
\text{固定候选参考law}\ \Longrightarrow\ \text{完整动作评分校准}\ \Longrightarrow\ \text{候选保留保证}.
\]
前者针对未知响应参数，后者针对已声明参考分布；它们不能未经验证互相替代。

\begin{center}\small
\begin{tabular}{L{3.0cm}L{11.6cm}}\toprule
记号 & 含义与维度\\\midrule
$O,\ n,\ T$ & 观测叶集合、叶数、每批槽数；$i,j\in O$，$t=1,\ldots,T$\\
$P,Q,Y$ & $T\times n$总表功率及平方电压降，$Y_{ti}=V_{0t}^2-V_{it}^2$；已包含EV实际注入\\
$\mathcal L,\ R,X$ & 观测后代支持的层状族及$n\times n$有功/无功响应矩阵；$C\in\mathcal L$为集合\\
$U,H,w$ & $T\times k$探测输入、$T\times d_H$干扰设计和$T$维对比权重\\
$\mathcal G_{r-1}$ & 第$r$批开始前的信息；不含本批计量噪声\\
$\mathcal A,\mathcal H,N$ & 冻结动作库、候选集合、每参考law的独立校准episode数\\
$B,\alpha,\delta,\varepsilon$ & 最多批次数、目标风险、校准坏事件概率、未来批越界率\\
\bottomrule\end{tabular}
\end{center}
功率和电压工作模型采用一致标幺值，$R/X$系数已包含平方电压模型的倍数2；设备调度另用kW、kvar、kWh，进入辨识模型前须换基值。
层状族指任意两支持不交或一者包含另一者。

\subsection{从径向线路到共享路径响应核}
假设网络是固定平衡径向树，根为0，$O$仅含观测叶；删去无观测后代分支，并将具有相同观测后代集合的串联边聚合。
所得\emph{观测缩减树}除物理根外的内部节点至少有两个观测后代子分支。
在$\mathcal L$中补入$O$和全部单叶支持，缺失时赋零权重；公共根干线聚合为$r_O,x_O$，物理根可只有一个观测子分支。
允许零权边，但零权分叉不承诺唯一可恢复。
隐藏支路上的源最多定位到对观测叶的投影接点；源自身无观测支线阻抗可能不可辨。

在线性无损支路模型中，支路功率及根到$i$的平方电压降满足
\[
 P_e=\sum_{j\in C_e}p_j,\qquad Q_e=\sum_{j\in C_e}q_j,\qquad
 Y_i=\sum_{e\in\operatorname{path}(0,i)}2(r_eP_e+x_eQ_e).
\]
交换求和并把同支持的$2r_e,2x_e$聚合为$r_C,x_C$，得到下列结构\cite{cavraro,taheri}。
\begin{proposition}[共享路径表示]\label{prop:kernel}
令$z_C=1_C$，$r_C,x_C\ge0$。则
\begin{equation}\label{eq:kernel}
 R=\sum_{C\in\mathcal L}r_Cz_Cz_C^\top,
 \quad X=\sum_{C\in\mathcal L}x_Cz_Cz_C^\top,
 \quad R_{ij}=\sum_{C\supseteq\{i,j\}}r_C.
\end{equation}
$R_{ij}$等于$i,j$共有根路径的有功权重和，且$R\succeq0$；
$d_{ij}=R_{ii}+R_{jj}-2R_{ij}$为叶间加性树距离，零权边存在时可为伪距离。
\end{proposition}
\begin{proof}
边$C$影响输出$i$且承载输入$j$恰当且仅当$i,j\in C$，故其矩阵贡献为$r_Cz_Cz_C^\top$。
对任意$v$，$v^\top Rv=\sum_Cr_C(z_C^\top v)^2\ge0$。
$d_{ij}$中共有根路径抵消，剩下从$i$经最近共同祖先到$j$的路径权重。
\end{proof}
半正定是必要性质，任意半正定矩阵不自动具有层状树分解。单列$R_{:,j}$只解析$j$祖先链上的划分，不能解析完全位于其他子树内部的分叉。
式\eqref{eq:kernel}是线性工作模型；非线性AC的损耗、工作点依赖和聚合误差必须进入下文$D$，不能因网络径向就令$D=0$。

\subsection{源项、误差集合与WZZT损失}
设$\widehat s_P,\widehat s_Q\in\mathbb R^T$为名义源幅值，$g_R,g_X\in\mathbb R_+^n$为该源对观测叶的共同路径响应。
工作模型为
\begin{equation}\label{eq:fullmodel}
Y=PR^\top+QX^\top+\widehat s_Pg_R^\top+\widehat s_Qg_X^\top+D+E.
\end{equation}
$D$包括源幅值偏差与AC近似误差，$E$为计量误差；等式成立本身不证明$D$小，必须独立给出覆盖真实允许域的界。
这里源项只表示未计入$P,Q$的注入，不能把已含在总表内的同一功率重复加入。

\begin{proposition}[共同源误差盒是外松弛]\label{prop:sourcebox}
固定全部响应参数，若$|\delta s_{Pt}|\le b_{Pt}$、$|\delta s_{Qt}|\le b_{Qt}$，定义
\[
\mathcal Z_t=\{g_Ru+g_Xv:|u|\le b_{Pt},|v|\le b_{Qt}\},\quad
 a_{ti}=b_{Pt}g_{Ri}+b_{Qt}g_{Xi}.
\]
则$\mathcal Z_t\subseteq\mathcal B_t:=\prod_i[-a_{ti},a_{ti}]$。对名义残差$e_t$，
\begin{equation}\label{eq:boxloss}
 \min_{d_t\in\mathcal B_t}\|e_t-d_t\|_1
 =\sum_i[|e_{ti}|-a_{ti}]_+
 \le\min_{d_t\in\mathcal Z_t}\|e_t-d_t\|_1.
\end{equation}
因此在同一参数域上继续最小化，盒损失最优值仍不高于共同源损失最优值。
\end{proposition}
\begin{proof}
三角不等式给$|g_{Ri}u+g_{Xi}v|\le a_{ti}$。
盒的各坐标可独立投影；标量$e$到$[-a,a]$的绝对距离为$[|e|-a]_+$。可行集合包含关系给最后一个不等式；逐参数成立的不等式在同域取下确界后仍成立。
\end{proof}
取$g_R=(1,1)^\top$、$b_P=.1$、$g_X=0$，残差$(.1,-.1)^\top$的盒损失为0，却不存在同一个源幅值消除它。
固定支持/源位置且$g$对阻抗为线性时，以$u_{ti}\ge0$、$u_{ti}\ge\pm e_{ti}-a_{ti}$构造LP，目标为$(Tn)^{-1}\sum_{ti}u_{ti}$。
这是盒问题的精确LP，仍不是一般共同源问题的精确等价。

若源投影在压缩边$C$内，则
\begin{align}\label{eq:sourceedge}
 g_R&=\sum_{A\supsetneq C}r_A1_A+r_C^{\rm pre}1_C,
 &0\le r_C^{\rm pre}\le r_C,\\
 g_X&=\sum_{A\supsetneq C}x_A1_A+x_C^{\rm pre}1_C,
 &0\le x_C^{\rm pre}\le x_C.
\end{align}
\textbf{双向构造证明：}任意合法前后串联段给出上述系数；反之，取前段$(r_C^{\rm pre},x_C^{\rm pre})$和后段$(r_C-r_C^{\rm pre},x_C-x_C^{\rm pre})$即构造一个等效串联表示。
所以对自由非负未知前后阻抗是精确重参数化；已知均匀导线则另需存在$\lambda\in[0,1]$使$(r_C^{\rm pre},x_C^{\rm pre})=\lambda(r_C,x_C)$，不能用自由分段放宽后还声称满足该物理约束。
路径表示还给出$0\le g_{Ri}\le R_{ii}$、$0\le g_{Xi}\le X_{ii}$，但工程上界不能从一个拟合解反推。
源位置域若允许根部，应包含$g_R=g_X=0$；其它接入位置、材料比例和正下界也须显式进入允许域。

\section{从有效输入到严格响应区间}
对一个输出，写成
\begin{equation}\label{eq:contrastmodel}
 y=U\beta+H\gamma+d+\epsilon.
\end{equation}
单向P且探测Q为0时$\beta=R_{ij}$；独立PQ时$\beta=(R_{ij},X_{ij})^\top$。
$H$收纳准确观测的背景功率、名义源及预定趋势，$\gamma$为未知系数，$d$容纳未消除偏差。

\subsection{何时存在无偏对比}
取$M_H=I-HH^\dagger$为$\operatorname{col}(H)$的正交补投影。
\begin{proposition}[全分量可估计的秩条件]\label{prop:rank}
对每个$l=1,\ldots,k$，存在$H^\top w_l=0$、$U^\top w_l=e_l$，当且仅当$\operatorname{rank}(M_HU)=k$。满秩时可取
\[
 w_l=M_HU(U^\top M_HU)^{-1}e_l.
\]
\end{proposition}
\begin{proof}
直接代入验证充分性。若不满秩，则存在$v\ne0$使$M_HUv=0$，即$Uv=Hz$。
若全部对比存在，则每个$v_l=w_l^\top Uv=w_l^\top Hz=0$，矛盾；同时$(\beta+v,\gamma-z)$与$(\beta,\gamma)$产生同一均值，说明在无附加参数限制的线性模型中完整$\beta$不可辨。
\end{proof}
单个分量只需$e_l$处于$U^\top\ker(H^\top)$中；不能把全秩写成每个单分量的必要条件。
不可辨构造还须满足实际参数域；有非负、离散或其它先验限制时，秩亏不自动说明某个受限边界点有多个可行解释。
单向P对应$M_Hp\ne0$；独立估计R/X需要两个残差化输入不共线。
若$Q=\kappa P$，只能估计$R+\kappa X$；其树支撑仍满足下文正权论证，须额外要求各有效边权$r_C+\kappa x_C\ge0$，目标边严格正。容性$\kappa<0$可抵消甚至翻转边权，不能自动沿用R列的clade证书。
例如$O=\{1,2,3\}$、$v=(1,1,0)^\top$，令$R=2\,11^\top+vv^\top+3I$、$X=2vv^\top$，两者的非平凡支持均只可能为$\{1,2\}$。
但$\kappa=-1$时$(R-X)_{:,1}=(4,1,2)^\top$，阈值1.5给出假支撑$\{1,3\}$；组合矩阵即使正定，也不能替代有效边权非负条件。

\subsection{偏差、噪声与浮点等式残差}
\begin{proposition}[条件覆盖与数值泄漏修正]\label{prop:contrast}
条件于不含本批计量噪声的历史、物理轨迹和设计信息。设计、权重、偏差界与协方差上界均已确定，且
\[
 |d|\le b,\quad \epsilon\sim N(0,\Sigma),\quad
 \Sigma\preceq\overline\Sigma,\quad
 |\beta|\le B_\beta,\quad |\gamma|\le B_\gamma.
\]
允许$w$不精确满足对比等式。设$\widetilde\beta_l$满足$|\widetilde\beta_l-w^\top y|\le e_{\rm eval}$，取
\begin{align}
 \rho_l={}&b^\top|w|+z\sqrt{w^\top\overline\Sigma w}
 +|U^\top w-e_l|^\top B_\beta
 +|H^\top w|^\top B_\gamma+e_{\rm eval}.\label{eq:v4contrast}
\end{align}
则$\Pr(|\widetilde\beta_l-\beta_l|>\rho_l\mid\text{上述信息})\le2\Phi(-z)$，其中$z\ge0$。
精确满足等式时两项泄漏为0，无需相应参数有限界。
\end{proposition}
\begin{proof}
展开$w^\top y-\beta_l=(U^\top w-e_l)^\top\beta+(H^\top w)^\top\gamma+w^\top d+w^\top\epsilon$。
前三项的绝对值分别由三个盒的支撑函数界住；最后一项条件高斯，方差不超过$w^\top\overline\Sigma w$。加入求值误差后应用双侧尾界。零方差时该随机项几乎处处为0。
\end{proof}
$w$不能根据本批带噪响应再择优；闭环控制使用当前噪声后，即使给定已实现设计，也不自动保留上述条件高斯分布。
若偏差界仅在事件$G_d$上成立，则失败事件包含于$G_d^c$与高斯尾事件之并，不需两者独立。
带噪PQ是误差变量问题，需要新增误差预算或另一套推断，不能直接套本命题。

幅值计量为$V_i+\eta_i$时，平方电压降噪声严格展开为
\[
 2V_0\eta_0-2V_i\eta_i+\eta_0^2-\eta_i^2.
\]
给定不依赖该噪声的真实轨迹，前两项是Gaussian线性变换；在$|\eta_i|\le u_i$的同时事件上，后两项绝对值不超过$u_0^2+u_i^2$，可并入$b$。
共同根误差允许输出相关；时间相关必须进入有效的协方差上界。真实轨迹随机时，先对每条轨迹给统一尾界，再积分。

\textbf{数值实现边界：}独立工具\path{numerical_guard_v4.py}对存储浮点的泄漏点积作精确有理运算并向上舍入，也核对最终线性估计的舍入误差。
它尚未接入历史批次区间/证书流水线，且不等于对协方差、平方根、分位数、物理量化及区间端点都做了形式化包围。
因此历史结果不能被追认为已自动使用式\eqref{eq:v4contrast}完整数值修正。若干扰系数无界，$10^{-10}$的等式残差本身也不能证明泄漏小。

\subsection{半宽优化为何有用}
精确等式下令$L^\top L=\overline\Sigma$，则最小半宽等价于
\begin{align*}
 \min_{w,v,s}\quad&b^\top v+zs,\\
 \text{s.t.}\quad&-v\le w\le v,\quad \|Lw\|_2\le s,\quad
 H^\top w=0,\quad U^\top w=e_l.
\end{align*}
这是SOCP，因为目标线性、约束为仿射或二阶锥；有限参数盒下，加入绝对值泄漏惩罚仍是凸问题。
任意事前可行权重配以有效半宽均有覆盖依据，求到最优只改善效率\cite{armstrong}。
现有实现用SLSQP寻找候选权重并以FWL回退，不将求解状态解释为SOCP全局最优证书。

\begin{proposition}[有界偏差无法靠重复采样消失]\label{prop:biasfloor}
若单向P满足$0\le p_t\le P_{\max}$、$P_{\max}>0$，$H$含截距，且每槽偏差可独立取$[-b,b]$中的任意值，则任何精确对比$w^\top p=1$、$w^\top1=0$的最坏偏差半宽至少为$2b/P_{\max}$。
\end{proposition}
\begin{proof}
$1=|w^\top(p-P_{\max}1/2)|\le(P_{\max}/2)\|w\|_1$，而盒上的最坏偏差恰为$\sup_{|d_t|\le b}|w^\top d|=b\|w\|_1$。
\end{proof}
若系数域允许相应平移，斜率相差$\Delta\beta$、截距相差$-\Delta\beta P_{\max}/2$的两个模型，均值差为$\Delta\beta(p_t-P_{\max}/2)$。
当$|\Delta\beta|\le4b/P_{\max}$时，两侧偏差盒可吸收该差异。这给出偏差地板的模型不可区分解释，且不依赖样本数。

\section{响应区间怎样认证树支撑}
以下将某节点的全部观测叶后代集合称为树支撑（clade）。
\begin{proposition}[单列clade与聚合间隔]\label{prop:clade}
在命题\ref{prop:kernel}的观测缩减树中，固定探测叶$j$，假设所有$i$同时满足$L_i\le R_{ij}\le U_i$。
对任意$j\in C\subsetneq O$、$|C|\ge2$，若
\begin{equation}\label{eq:v4clade}
 \Delta_C:=\min_{i\in C}L_i-\max_{i\notin C}U_i>0,
\end{equation}
则$C\in\mathcal L$且$r_C\ge\Delta_C$。
\end{proposition}
\begin{proof}
在式\eqref{eq:v4clade}两端之间选阈值$c$。由同时覆盖，$C=\{i:R_{ij}>c\}$。
共有路径深度随$j$的祖先链单调不减，严格上水平集合必是这条链上某节点的全部观测后代；零权边只会合并层次。因此$C$是一个clade。
记$\pi(C)$为最近严格父支持。缩减树内没有单观测子分支：$C$中存在与$j$分属两个直接子分支的叶，$\pi(C)\setminus C$也有叶。因而
\[
 \min_{i\in C}R_{ij}=d_R(C),\qquad
 \max_{i\notin C}R_{ij}=d_R(\pi(C)),\qquad
 r_C=d_R(C)-d_R(\pi(C))\ge\Delta_C.
\]
其中$d_R(C)$为根到支持$C$的聚合深度。
\end{proof}
只要区间已同时覆盖，就可事后搜索全部$C$而不再按集合数扣减风险，因为覆盖事件上\emph{每一个}通过的证书均正确。
证书只下界聚合边权，不能把该权重分配给无法观测的串联物理边。X列同理；混合响应还须检查有效边权非负。

\begin{proposition}[恢复的充分条件与停止边界]\label{prop:stop}
固定非平凡目标族$\mathcal L_\star\subseteq\{C\in\mathcal L:2\le|C|<n\}$。若每个$C\in\mathcal L_\star$至少被一次$j\in C$的探测覆盖，所有对应区间为$[\widehat R_{ij}-\rho_i,\widehat R_{ij}+\rho_i]$且$\rho_i\le\rho$，全部目标聚合权重满足$r_C>4\rho$，则在同时覆盖事件上，按该列估计中心从大到小排序、遍历全部相邻分割并检验式\eqref{eq:v4clade}，必能认证所有目标clade。
\end{proposition}
\begin{proof}
覆盖给出$L_i\ge R_{ij}-2\rho$、$U_i\le R_{ij}+2\rho$；代入上一命题中的真实间隔，得$\Delta_C\ge r_C-4\rho>0$。
于是每个内部估计中心大于每个外部估计中心，$C$必为该估计排序的一个前缀，实际枚举会访问它；对每个目标选其已探测到的叶即可。
\end{proof}
这是对输入/位置覆盖、正间隔及区间宽度的充分条件，未证明当前WZZT启发式一定访问所需节点或穷尽全部目标。
错误认证控制与最终停止是两个命题。只有偏差半宽低于间隔、剩余噪声宽度能缩小且资源允许时才可能停止；有界但任意的逐槽偏差不会自动按$T^{-1/2}$消失。
$n-2$个不同且相容的非平凡真clade仅在目标确为完整二叉观测层级时给出全部内部层级；不说明全部参数或隐藏物理节点已恢复。

\section{三套置信方法：分别对什么概率负责}
\subsection{A：条件区间、风险分配与可选停止}
第$r$批行动由$\mathcal G_{r-1}$确定；采集新批且条件噪声/偏差假设成立。令$B_r$为该批任一所用区间失败，分配事前或可预测风险$\alpha_r$使
\[
 \Pr(B_r\mid\mathcal G_{r-1})\le\alpha_r,\qquad
 \sum_{r\ge1}\alpha_r\le\alpha\quad\text{几乎处处}.
\]
例如有限$B$批、每批$n$输出及$k$输入方向可取$z=\Phi^{-1}(1-\alpha_L/(2nkB))$；平方计量的有界事件再分配总风险$\alpha_Q$。
\begin{proposition}[跨批误认证控制]\label{prop:spending}
上述条件下，任意停止时刻之前曾输出错误clade的概率至多$\alpha$。
\end{proposition}
\begin{proof}
所有通过区间认证的clade在区间覆盖时都正确，故误认证事件包含于$\bigcup_r B_r$。并集界、塔式期望及单调收敛给
$\Pr(\bigcup_rB_r)\le\mathbb E\sum_r\Pr(B_r\mid\mathcal G_{r-1})\le\mathbb E\sum_r\alpha_r\le\alpha$。
可选停止只取其子事件，不增加该界。
\end{proof}
因此合法风险分配本身也能支持可选停止，无需将每轮固定95\%区间无限复用；理论的无限轮结论不代表当前有限轮实现已自动采用可求和预算。

\subsection{B：逐动作PAC包络}
给定独立训练信息，冻结$A=|\mathcal A|$个动作，每动作包括节点、时窗、完整PQ及相关初始状态。
对每动作$a$，$S_{1a},\ldots,S_{Na}$为其\emph{独立同分布}参考批的整批最大误差评分，分布函数$F_a$固定；不同动作之间允许相关。
\begin{proposition}[非参数容忍上界]\label{prop:pac}
令$q_a=S_{(k),a}$。若
\begin{equation}\label{eq:pactail}
 \Pr\{\operatorname{Bin}(N,1-\varepsilon)\ge k\}\le\delta/A,
\end{equation}
则以校准概率至少$1-\delta$，所有动作同时满足$\Pr_{\rm new,a}(S>q_a)\le\varepsilon$。
\end{proposition}
\begin{proof}
用随机化概率积分变换耦合出独立$U_b\sim\mathrm{Unif}(0,1)$，满足$U_b\le F_a(S_{ba})$，从而$U_{(k)}\le F_a(S_{(k),a})$。
坏事件$1-F_a(q_a)>\varepsilon$蕴含$U_{(k)}<1-\varepsilon$；其概率为或不超过式\eqref{eq:pactail}的二项尾。最后对$A$动作并集。分数有原子/并列也成立，只会保守\cite{wilks,vovk}。
\end{proof}
存在可用$k\le N$的最小样本要求为$N\ge\lceil\log(\delta/A)/\log(1-\varepsilon)\rceil$，因为最小尾概率在$k=N$时为$(1-\varepsilon)^N$。不足时取$q_a=+\infty$。
在校准好事件上，若给定校准集和每批过去，所选动作的新批仍服从该固定参考law，则至多$B$批的错误证书概率有界于
\[
 \min\{1,\ \alpha_V+\delta+B\varepsilon\},\qquad \alpha_V=\alpha_L+\alpha_Q.
\]
证明为校准坏事件、每批包络越界和计量区间失败三者并集。仅有可交换性不足以推出此条件于已实现校准集的PAC结论；实际计算二项尾仍有浮点精度边界。

\subsection{C：整个episode的动作最大值与候选反演}
条件于独立训练集$\mathcal D_0$，冻结有限候选$\mathcal H$、动作库$\mathcal A$和完整评分程序$s_h$。
对每个$h$，必须从\emph{$h$自身为真}的参考分布生成潜在结果向量$Z_b=(Y_b(a):a\in\mathcal A)$，使校准$Z_1,\ldots,Z_N$与一个新$Z_{N+1}$可交换。$Z_b$可含相关时间点、多个动作及接入日程；独立或可交换的单位是整个episode。

\begin{proposition}[同时动作下的边际候选覆盖]\label{prop:maxconformal}
定义$S_{bh}=\max_{a\in\mathcal A}s_h(Y_b(a),a)$，
\[
 k=\lceil(N+1)(1-\alpha)\rceil,\quad
 q_h=\begin{cases}S_{(k),h},&k\le N,\\+\infty,&k>N.\end{cases}
\]
新episode实际观察任意子集$K\subseteq\mathcal A$，取
\[
 \mathcal C=\{h:\max_{a\in K}s_h(Y(a),a)\le q_h\},
\]
并约定$K=\varnothing$保留全部候选。若真实参考候选为$h_\star$，则
\begin{equation}\label{eq:candidatecoverage}
 \Pr_{h_\star}(h_\star\notin\mathcal C\mid\mathcal D_0)\le\alpha,
 \qquad
 \Pr_{h_\star}(\mathcal C=\{h\},\ h\ne h_\star\mid\mathcal D_0)\le\alpha.
\end{equation}
\end{proposition}
\begin{proof}
给并列分数添加独立随机次序后，新$S_{N+1,h_\star}$的秩在$1,\ldots,N+1$均匀；事件$S_{N+1,h_\star}>S_{(k),h_\star}$概率至多$(N+1-k)/(N+1)\le\alpha$。
非严格保留使并列只增加保守性。对每条实现的潜在结果，子集最大值不大于全库最大值，所以真候选误排包含于上述越界事件。
错误单元素集合必然排除$h_\star$，给出第二个界；无需再按候选数分配$\alpha$。
\end{proof}
$K$可由本episode中已观测数据决定，因为上述包含关系是逐实现成立的；但真实动作必须等于校准潜在结果中的对应分量。
若执行先前动作改变后续SOC、工作点或控制规则，应将整个可执行历史策略作为新的episode定义，不能直接套独立初态动作库。
失败、空集与多候选不应被强制变成一个拓扑。模板不含真实参数/分布时，式\eqref{eq:candidatecoverage}不适用。

这是给定训练信息后的\emph{校准集与新episode共同随机的边际结论}，不是给定已经实现的校准集、天气或“已输出单例”的条件结论。
若需以$1-\delta$校准概率同时约束各候选参考law的未来误排率，可对$S_{bh}$应用命题\ref{prop:pac}，将式\eqref{eq:pactail}右端换成$\delta/|\mathcal H|$；这里只在更强校准陈述中支付候选数的$\delta$预算。

\subsection{相对评分的原理及不能跨越的概率边界}
已有相对评分为$m_h=s_h-\min_{g\ne h}s_g$，再用固定$\operatorname{expit}$映射。
对任意所有候选共有的$c(Y,a)$，以$s_h+c$替换$s_h$时$m_h$完全不变。这解释其抵消共同残差、突出结构竞争的机制；同时会丢失所有模型都很差的信息。
命题\ref{prop:maxconformal}只使用完整评分程序固定及交换性，故不要求它是精确似然比。方法换掉后必须用新独立校准或已预先覆盖的方法轴；不允许看到当前结果后择取最有利方法。

令$A$表示“输出单元素集合”，$W$表示错误输出。边际保证只给
\[
 \Pr(W)\le\alpha,\qquad
 \Pr(W\mid A)=\frac{\Pr(W)}{\Pr(A)}\le\min\{1,\alpha/\Pr(A)\},\quad\Pr(A)>0.
\]
当输出罕见时，右式可无信息；例如以$\alpha$概率输出错误单例，其余输出全候选，仍可有$1-\alpha$覆盖而每次单例都错。
因此覆盖、错误认证、选择后条件错误、最终停止与参数精度须分别陈述。

\section{为什么选择波形、地点与时段}
\subsection{两固定分布的可分性基准}
\begin{proposition}[Gaussian强制分类基准]\label{prop:gaussian}
若两候选等先验、$Y\mid h\sim N(\mu_h,\Sigma)$、共同已知$\Sigma\succ0$，定义$d^2=(\mu_1-\mu_0)^\top\Sigma^{-1}(\mu_1-\mu_0)$。必须二选一的分类器之最小平均错误率恰为$\Phi(-d/2)$。
\end{proposition}
\begin{proof}
对数似然比为$(\mu_1-\mu_0)^\top\Sigma^{-1}(Y-(\mu_1+\mu_0)/2)$，两假设下分别为$N(\pm d^2/2,d^2)$。
逐样本选择较大密度最小化等先验平均损失，零阈值给两侧错误$\Phi(-d/2)$。$d=0$时分布相同，最小平均错误为$1/2$。
\end{proof}
达到5\%强制分类错误至少要求$d\ge2\Phi^{-1}(.95)\approx3.29$；它不直接约束允许弃权的唯一认证成功率。
白化后只观察正交补投影时，距离为
\[
 d_\perp^2=\|(I-\Pi_{\Sigma^{-1/2}H})\Sigma^{-1/2}(\mu_1-\mu_0)\|_2^2\le d^2.
\]
取正交补的一组标准正交基可将奇异投影写成非退化Gaussian观测，重复上一证明；这只是该投影实验的最优错误，非任意复合模型GLRT定理。

\subsection{单向充电可提供多少激励}
端口$j$额外EV引起的线性均值增量差为$\Delta R_{:,j}p_t+\Delta X_{:,j}q_t$；若这两列均无差异，增加该方向激励不能区分这对固定候选。
固定波形下，地点改变响应差，时段改变干扰/噪声及可接入性；这些因素必须来自执行前的信息。

\begin{proposition}[固定充电量的离散激励上界]\label{prop:energy}
设$P_{\max}>0$、$0\le p_t\le P_{\max}$、$K_p=\sum_{t=1}^Tp_t$固定，$H$包含常量列。则
\[
 \|M_Hp\|_2^2\le\sum_t(p_t-K_p/T)^2
 \le P_{\max}K_p-K_p^2/T.
\]
更紧的离散上界为$mP_{\max}^2+\zeta^2-K_p^2/T$，其中$m=\lfloor K_p/P_{\max}\rfloor$、$\zeta=K_p-mP_{\max}$。
\end{proposition}
\begin{proof}
正交投影到更小子空间不增加范数；再用$p_t^2\le P_{\max}p_t$求和。
若两个分量都严格在$(0,P_{\max})$内，把较小分量的功率移向较大分量直到一个到端点，会增加平方和且保持总量，故最大平方和在至多一个非端点的向量上取得，给出离散上界。
\end{proof}
理想0/$P_{\max}$占空码在$K_p/P_{\max}$为整数时达到第二个界；若还与中心化时间正交，则消去线性趋势不再损失该范数。
常量充电被截距消去。延迟、爬坡、SOC及驻留约束会缩小可行域，所以上界不是可执行性保证。
本例$T=24$、$K_p=9/0.25=36$kW、$P_{\max}=6$kW，理论上界162kW$^2$；已测稀疏码133.76、0/3kW码33.97，收益同时包含幅值和占空比变化。
此量是欧氏激励能量；只有标量独立同方差Gaussian回归时，除以$\sigma^2$才是消干扰后的Fisher信息，不能无条件当拓扑置信度。

\subsection{优化目标与下界方向}
\begin{proposition}[局部注入对全方向激励的限制]\label{prop:lowrankprobe}
对固定同一时窗及预处理下的$S\in\mathbb R^{d\times T}$，定义$s_i(S)=\sqrt{\lambda_i(SS^\top)}$，按非降序排列，故$T<d$时自动含零。
若EV新增输入$E$满足$\operatorname{rank}(E)\le r<d$，则
\[
 s_1(S+E)\le s_{r+1}(S).
\]
\end{proposition}
\begin{proof}
取$SS^\top$最小$r+1$个特征值对应向量张成空间$V$。
由$\dim V+\dim\ker(E^\top)\ge d+1$，存在单位向量$v\in V\cap\ker(E^\top)$，于是
\[
 s_1(S+E)\le\|(S+E)^\top v\|_2=\|S^\top v\|_2\le s_{r+1}(S).
\]
\end{proof}
若同一时段的背景不随EV响应，则单节点P或独立PQ输入分别至多增加1或2秩；功率再大，也不能保证全部薄弱方向被改善。
新增观测时长、变换预处理或背景反馈须重新建模，不能直接套该比较。
这不排除局部辨识：例如$S=(1,0)^\top$的全方向最小奇异值为0，但若允许的竞争响应差仅为$\Delta A=(\delta,0)$、$\delta\ge\delta_0>0$，则$\|\Delta A S\|\ge\delta_0$。
因此可以研究争议域内的受限增益或响应列区间；其有效性仍需覆盖完整竞争域，不能由固定模板的成功率倒推。

给定工程上可执行目录\cite{bhela}，D/E设计优化统一参数维度的信息矩阵行列式/最小特征值，T型设计针对最难候选对的响应差\cite{dette}。
前者可能在相同输入和噪声的不同端口上打平；后者才直接针对当前结构争议。未知参数时可研究
\begin{equation}\label{eq:robustdesign}
 \max_{a\in\mathcal A_{\rm feasible}}\frac{1}{c_a}
 \min_{h\ne g}\inf_{\theta\in\Theta_h,\theta'\in\Theta_g}
 \|\Sigma_a^{-1/2}(\mu_h(a,\theta)-\mu_g(a,\theta'))\|_2^2,
 \qquad c_a>0.
\end{equation}
内层可行局部解或有限采样给最小值的\emph{上界}，不能证明统一正分离；需要可靠\emph{下界}。两个允许域存在同分布解释时，该方向最坏分离为0。
式\eqref{eq:robustdesign}仍是设计目标，不自动构成统计检验；已有主试验实现固定模板评分，未实现全未知参数双侧全局优化。
算法输出应包含接入点、窗口、可执行PQ、目标结构争议、服务/作业成本及未决原因。优化用于减少采样成本和区间宽度，最终认证仍取决于新数据、有效界与明确概率目标。


\subsection{与WZZT衔接的执行流程及算法范围}
原完整边源模型为$g_R=\sum_k r_kh_kz_k$、$g_X=\sum_k x_kh_kz_k$，其中$h$是一条完整祖先路径。
每次支持扩展、重拟合、剪枝与交换都重新优化整个当前$h$，不能将旧$h$冻结后只加一项。
压缩边前/后阻抗模块扩充的是源位置表示；它不自动把共同源幅值误差的外松弛变为精确模型。
\begin{enumerate}
\item 被动数据产生暂定WZZT支持族和候选；将已独立认证支撑与候选提示分开。
\item 对尚未认证的支撑生成竞争解释，冻结设备、时间和接入动作库。地点评分读取过去拟合或工程模板，当前Y尚不可用。
\item 依据目标选择局部响应D/E、候选T区分或最小置信半宽；通过接入、设备及可用网络约束后执行一个动作，记录实际功率。
\item 用新数据构造响应证书或候选集合。无覆盖条件、校准不足、空集或多候选时保留未决。
\item 在新数据上重估全部R/X及完整源路径；已认证支持及下界受保护，未认证支持允许删除/交换。结构满足要求后，参数精度另行检查。
\end{enumerate}
现有WZZT的一支持扩展只证明固定已有支撑下的单次最优扩展；局部删除/交换不等于全树全局最优。
v3已实现证书账本和多轮结构修正；v4主试验单独检验有限工程模板下的时地选择，未知系数衔接由后文小算例验证。没有把二者拼成已验证的任意网络全流程保证。

\subsection{证明和实现的对应关系}
\begin{center}\small
\begin{tabular}{L{4.0cm}L{5.6cm}L{5.0cm}}\toprule
命题/机制 & 已实现证据 & 尚不能声称\\\midrule
源盒LP与完整源路径重拟合 & 固定族/单次扩展；边内源枚举 & 所有未知树的全局最优\\
条件响应区间与clade & 有界偏差、FWL/优化权重、有限轮证书 & 任意带噪PQ、任意源域覆盖\\
式\eqref{eq:v4contrast}泄漏修正 & 独立精确点积工具及单测 & 已自动接入历史区间流水线\\
命题\ref{prop:maxconformal} & 固定候选、完整60动作、两种独立校准评分 & 全未知阻抗复合域或任意SOC继承路径\\
式\eqref{eq:lpbound}/\eqref{eq:rhsbound} & 校正标量下界；原始局部价格 & 修正价格与标量界已成套导出\\
连续EV设备模型 & 轨迹生成与浮点容差审计 & 实车跟踪或全过程网络安全证明\\
命题\ref{prop:eprocess}/\ref{prop:selective} & 本文条件推导 & 已实现的序贯/选择性风险保证\\
\bottomrule\end{tabular}
\end{center}
数学证明使用精确实数及明确概率假设；代码核查与历史数值试验是另一层证据。两者之间的接口缺口按表保留。


\clearpage
\section{组织算例：先冻结设计，再检验收益}
\subsection{三个层次的试验}
\begin{center}\small
\begin{tabular}{L{2.3cm}L{6.5cm}L{5.8cm}}\toprule
层次 & 数据与比较 & 能回答的问题\\\midrule
结构/置信基线 & 既有五末端v3：工程响应带、PAC包络、全源路径拟合 & 条件clade证书与未知系数拟合的数学机制\\
时地效果主研究 & 十末端、三固定阻抗模板、60动作、13策略，真实平衡AC & 控制类型、地点与时段选择对已声明模板的判别收益\\
设备/误差补充 & 最小驻留/延迟、稀疏单向码、细步AC积分、源参数重拟合 & 可执行假设是否影响收益，哪些结论不能推广\\
\bottomrule\end{tabular}
\end{center}
主研究的三个模板只改变叶5--8的二二配对：$(56|78)$、$(57|68)$、$(58|67)$。
网络有10个观测末端和8个根/内部母线。接入点1/9对这些模板无区分，5/7有区分；这是可检查的结构机制，不是假设算法已知真模板。
候选R/X由三个公开工程模板给定；选点可使用全部模板，但不读取哪一个是真值。
因此主结果是\textbf{固定模板识别}，不能用作任意未知阻抗完整WZZT恢复率。

\subsection{地点、时段和公平的控制组}
每次派一辆车、一个地点、一个6h时窗；潜在动作库仅供校准和选择，不是让每个策略实际执行60次。
动作包含无车、恒功率、单向P、单向PQ、双向PQ，乘4地点和3时窗。
接入日程在量测前确定；不可接入返回未决，并保留在分母。
本版只给事前接入日程与固定驾驶能耗，没有地理路网、到达时间或员工班次优化；地点收益来自避开候选盲端口及适应可接入性，尚不能换算为现场车队节省的工时。
\begin{center}\small
\begin{tabular}{lrrr}\toprule
窗口 & 背景负荷倍率 & 隐藏源随机波动尺度(pu) & 电压噪声标准差(pu)\\\midrule
00--06 & 0.65 & 0.0004 & 0.00010\\
10--16 & 1.00 & 0.0020 & 0.00024\\
18--24 & 1.45 & 0.0060 & 0.00055\\
\bottomrule\end{tabular}
\end{center}
这是明确的日内合成参考分布，噪声随时段变化代表事前监测信息，尚非现场证据。
若信息质量、可接入性及背景完全相同，时段优化没有必然优势。
低压为400V/100kVA；中压为11kV/1MVA、实际阻抗为低压的75.625倍、背景kW为10倍，同一车的标幺注入缩小10倍。
中压车载PQ应解释为\textbf{经理想配变折算的等效注入}；不表示EV直接接11kV，未计配变/低压支线损耗及相别。

主研究每尺度、每真模板用30个独立训练episode冻结噪声评分尺度，400个新episode校准，200个新episode测试。
每个episode包含相关的完整60动作潜在结果，独立单位是seed而非槽或动作。
三个模板复用同seed随机流，所以600条模板记录仍只提供200个独立簇。
训练/校准/测试seed域分别从20000/40000/80000开始；正式测试之后未根据效果改变主参数。

得分是候选预测$PR_h^\top+QX_h^\top$的标准化残差，经功率首端代理、固定截距/趋势投影，并对末端中心化消去共同根电压量测。
归一化与拟合过程本身纳入参考校准。主研究的响应RMSE只针对\textbf{中心化响应列}，不能写成绝对R/X精度。
所有失败、不可接入、空集合及多候选集合保留；仅单元素置信集合输出唯一模板。
这里的保证是重复试验中的覆盖/错误认证概率，\textbf{不是观测到输出以后该拓扑有95\%后验概率}。

\section{EV与地点--时段选择的实际效果}
表中均为百分比；每尺度分母为三模板各200次、合计600次。覆盖是真模板在集合内，错误是错误单元素输出；覆盖率本身不要求唯一。
\begin{center}\footnotesize
\begin{tabular}{L{4.7cm}rrrrrr}\toprule
 & \multicolumn{3}{c}{低压参考域} & \multicolumn{3}{c}{中压折算参考域}\\
策略 & 唯一正确 & 覆盖 & 错误 & 唯一正确 & 覆盖 & 错误\\\midrule
无 EV 观察基线 & 0.0 & 98.5 & 0.0 & 0.0 & 98.0 & 0.0 \\
固定地点/晚间恒功率 & 0.0 & 99.0 & 0.0 & 0.0 & 98.5 & 0.0 \\
固定地点/晚间 P 调制 & 0.7 & 98.5 & 0.0 & 0.0 & 98.5 & 0.2 \\
固定地点/晚间 PQ & 11.2 & 98.5 & 0.0 & 0.0 & 98.5 & 0.0 \\
固定地点/晚间双向 PQ & 10.3 & 98.5 & 0.0 & 0.0 & 98.5 & 0.2 \\
晚间随机可用地点 PQ & 6.5 & 98.0 & 0.0 & 0.0 & 98.0 & 0.0 \\
晚间 T 型选点 PQ & 12.3 & 98.0 & 0.0 & 0.0 & 98.0 & 0.0 \\
固定地点择时 PQ & 92.8 & 98.5 & 0.0 & 4.8 & 98.0 & 0.0 \\
T 型地点/时段联合 PQ & 97.0 & 97.5 & 0.0 & 6.0 & 96.0 & 0.0 \\
T 型地点/时段联合 P & 72.3 & 97.5 & 0.0 & 0.5 & 96.0 & 0.7 \\
T 型地点/时段联合双向 PQ & 96.0 & 96.5 & 0.0 & 18.3 & 97.0 & 0.0 \\
固定盲端口 PQ & 0.0 & 98.0 & 0.0 & 0.0 & 98.0 & 0.0 \\
局部列 D 信息择时/选点 & 0.0 & 97.5 & 0.0 & 0.0 & 96.0 & 0.0 \\
\bottomrule\end{tabular}
\end{center}
T型分数用最小候选响应距离与训练所得时段尺度排序，不能读当前Y；随机策略使用同一可用地点集合。
局部列D目标在相同输入/噪声的节点之间相同，表中按固定规则解平局，无法自动发现盲端口；这不否定针对全部参数的其他D设计。
双向PQ额外增加调制方差与循环量；对照只针对冻结码，不能推出V2G优于所有单向最优波形。
低压联合PQ相对只择时增加4.17个百分点，但全族Holm检验不显著；低压双向相对单向PQ也不显著。
中压双向PQ比单向PQ增加12.33个百分点，虽显著，仍有多数未决。中压P联合策略7/600次输出唯一候选，其中3次正确、4次错误：无条件错误输出率0.67\%，而输出条件下样本错误率为4/7。前者受覆盖理论约束，后者不能据此称为至多5\%。

\subsection{派车输出应包含什么}
例如低压seed 80000的联合策略在观察当前电压之前选\textbf{端口5、00--06、单向独立PQ}，目标争议为叶5--8的三种配对。成本为占用6h、网侧9kWh及驾驶/换接0.55kWh；实际电池增能8.358kWh、离站SOC约0.430。该例新数据仅保留模板0；单次成功只是完整测试中的一个记录。
下图给该动作24个实际槽均值；连线用于阅读，不代表瞬时阶跃。执行应调用已审计连续轨迹，不能直接把均值当成待执行指令。
\begin{center}
\begin{tikzpicture}
\begin{axis}[width=.46\textwidth,height=4.6cm,xlabel={时间(h)},ylabel={平均P(kW)},xmin=0,xmax=6,ymin=0,ymax=3.3,grid=major,tick label style={font=\small},label style={font=\small}]
\addplot[blue!65!black,thick,mark=*,mark size=1.1pt] coordinates {(0.125,2.49712) (0.375,0.34994) (0.625,0.07717) (0.875,2.66502) (1.125,0.34994) (1.375,2.66502) (1.625,2.93779) (1.875,0.34994) (2.125,0.07717) (2.375,2.66502) (2.625,2.93779) (2.875,0.34994) (3.125,0.07717) (3.375,2.66502) (3.625,2.93779) (3.875,0.34994) (4.125,2.07960) (4.375,0.34994) (4.625,0.07717) (4.875,2.66502) (5.125,0.93535) (5.375,2.66502) (5.625,2.93779) (5.875,0.33836)};
\end{axis}
\end{tikzpicture}\hfill
\begin{tikzpicture}
\begin{axis}[width=.46\textwidth,height=4.6cm,xlabel={时间(h)},ylabel={平均Q(kvar)},xmin=0,xmax=6,ymin=-5,ymax=5,grid=major,tick label style={font=\small},label style={font=\small}]
\addplot[orange!80!black,thick,mark=*,mark size=1.1pt] coordinates {(0.125,-3.82500) (0.375,-4.50000) (0.625,1.80000) (0.875,4.50000) (1.125,4.50000) (1.375,4.50000) (1.625,-1.80000) (1.875,-4.50000) (2.125,-4.50000) (2.375,-4.50000) (2.625,1.80000) (2.875,4.50000) (3.125,-1.80000) (3.375,-4.50000) (3.625,1.80000) (3.875,4.50000) (4.125,-1.80000) (4.375,-4.50000) (4.625,1.80000) (4.875,4.50000) (5.125,-1.80000) (5.375,-4.50000) (5.625,1.80000) (5.875,3.82500)};
\end{axis}
\end{tikzpicture}
\end{center}
对其他任务，派车记录还应给出预测最弱分离、适用误差域和未决原因。盲端口、功率不足、阻抗域重叠或无合法接入时，应返回相应原因及未消除的支撑争议，避免仅输出残差最大的节点。

\subsection{显著性：保留配对和多重比较}
差值定义为后者减前者，单位百分点。95\%区间是按200个完整seed簇重采样5000次的描述性bootstrap区间。
精确簇sign检验考察正/负改进簇概率，\emph{不直接等同于总体平均差为零的检验}；22项主比较跨两尺度统一Holm校正。
逐模板精确McNemar检验另保存在结果中，不能把三个相关模板当1800个独立检验样本。
\begin{center}\footnotesize
\begin{tabular}{L{5.1cm}crrr}\toprule
配对比较 & 域 & 差值 & 描述性95\%区间 & Holm $p$\\\midrule
同能量：恒功率→P调制 & LV & +0.7 & [0.2, 1.3] & $1$ \\
 & MV & +0.0 & [0.0, 0.0] & $1$ \\
同地点时段：P→PQ & LV & +10.5 & [7.3, 13.8] & $3.09\times10^{-11}$ \\
 & MV & +0.0 & [0.0, 0.0] & $1$ \\
同晚间：随机→T选点 & LV & +5.8 & [3.2, 8.8] & $5.91\times10^{-4}$ \\
 & MV & +0.0 & [0.0, 0.0] & $1$ \\
T选点：固定晚间→联合择时 & LV & +84.7 & [80.0, 89.0] & $1.19\times10^{-49}$ \\
 & MV & +6.0 & [3.8, 8.5] & $2.38\times10^{-7}$ \\
择时：固定地点→联合选点 & LV & +4.2 & [1.0, 7.7] & $0.2471$ \\
 & MV & +1.2 & [0.3, 2.2] & $0.3125$ \\
联合设计：单向PQ→双向PQ & LV & -1.0 & [-2.5, 0.0] & $1$ \\
 & MV & +12.3 & [9.0, 15.8] & $1.10\times10^{-13}$ \\
\bottomrule\end{tabular}
\end{center}
不显著只能说明本样本未分开相应策略，不能宣布等效；bootstrap用于效果评价，不负责拓扑认证。

\clearpage
\section{多套置信方法在相同数据上的代价}
既有路线A/B在五末端v3中已完成200个多轮测试：200次完成观测层级，零错误证书。
其风险为$\alpha_V+\delta+3\varepsilon=0.015+0.005+3(0.01)=0.05$，依赖10个冻结动作各1000个参考批次及相应历史条件分布。
这组结果保留为独立研究基线，不与十末端固定模板率拼接成一个恢复率。

对当前十末端同一批数据比较C路线的两个版本：split的单episode边际风险0.05；PAC取未来误排0.04和校准坏事件0.01，使单episode无条件上界同为0.05。
后者以至少99\%校准概率同时保证三个模板参考law各有至少96\%保留概率，代价是更高的次序阈值。
\begin{center}\small
\begin{tabular}{llrrrrr}\toprule
域 & 校准 & 次序$k$ & 唯一正确 & 真模板保留 & 错误唯一 & 平均集合大小\\\midrule
LV & split & 381 & 582/600 & 585/600 & 0 & 0.985 \\
LV & PAC & 395 & 595/600 & 600/600 & 0 & 1.015 \\
MV & split & 381 & 36/600 & 576/600 & 0 & 2.707 \\
MV & PAC & 395 & 25/600 & 597/600 & 1 & 2.863 \\
\bottomrule\end{tabular}
\end{center}
表中使用同一联合PQ策略；输出更宽集合是有意义的未决，不是应由最小残差强行补齐的错误。
这里的5\%是\textbf{一个新episode的风险}，不是在反复复用校准集的全部200次测试中错误总数为零的保证。
校准二项尾由SciPy计算并保存判据余量；统计定理成立的条件与浮点尾计算的精度分开表述。
原总残差评分在中压联合PQ下可置信唯一率只有6\%，但同数据强制选最小残差的正确率为92.5\%。这表明当前置信评分存在功效代价，6\%不能被误解为任何算法的可辨识性上限。

\subsection{方案C的改进：校准相对竞争分数}
在原研究揭示上述功效问题后，先固定新的完整评分程序，再采集全新的校准和测试模拟数据：
\[
 m_h(Y,a)=s_h(Y,a)-\min_{g\ne h}s_g(Y,a),\qquad
 \widetilde s_h=\operatorname{expit}(m_h).
\]
对所有候选共有的加性分数偏移$m_h$不变；它衡量本候选比最佳对手差多少。RMS归一、源代理重拟合与AC误差使它不能被称为精确似然比。
固定expit只为复用非负评分接口；校准直接作用于实际浮点评分，饱和并列以非严格阈值保留。
覆盖证明仍是方案C的交换秩证明，既不要求原分数最优，也不要求残差服从Gaussian。文献支持评分函数会影响集合效率\cite{margin}，并未证明margin对所有问题最优。

补充试验仅在中压参考域，沿用原30组训练所得的固定设计及完整60动作；每真模板用200个新校准、100个新测试episode，seed域430000/440000。三模板共享seed，独立测试单位为100个簇。两种评分分别按$\alpha=.05$校准；以下对比使用同批全新数据，不能将其绝对评分5.0\%与旧研究6.0\%混为同一样本。
\begin{center}\small
\begin{tabular}{L{4cm}lrrrr}\toprule
中压策略 & 评分 & 唯一正确\% & 覆盖\% & 错误\% & 平均集合大小\\\midrule
T 型地点/时段联合 P & 总残差 & 0.7 & 95.0 & 0.3 & 2.84 \\
 & 竞争差 & 6.3 & 99.7 & 0.0 & 2.62 \\
T 型地点/时段联合 PQ & 总残差 & 5.0 & 95.0 & 0.0 & 2.72 \\
 & 竞争差 & 74.7 & 99.0 & 0.7 & 1.36 \\
T 型地点/时段联合双向 PQ & 总残差 & 12.7 & 95.0 & 0.0 & 2.47 \\
 & 竞争差 & 88.7 & 99.7 & 0.3 & 1.15 \\
固定地点/晚间 PQ & 总残差 & 0.0 & 100.0 & 0.0 & 3.00 \\
 & 竞争差 & 0.0 & 100.0 & 0.0 & 3.00 \\
\bottomrule\end{tabular}
\end{center}
\begin{center}\footnotesize
\begin{tabular}{L{5cm}rrr}\toprule
总残差改为竞争差 & 差值(pp) & 描述性95\%区间 & 四比较Holm $p$\\\midrule
T 型地点/时段联合 P & +5.7 & [3.0, 8.3] & $1.53\times10^{-4}$ \\
T 型地点/时段联合 PQ & +69.7 & [62.7, 76.3] & $5.17\times10^{-26}$ \\
T 型地点/时段联合双向 PQ & +76.0 & [68.7, 83.0] & $7.75\times10^{-26}$ \\
固定地点/晚间 PQ & +0.0 & [0.0, 0.0] & $1$ \\
\bottomrule\end{tabular}
\end{center}
单向PQ相对评分共2/300次错误唯一，双向PQ共1/300次；这不是零风险方法。相对评分改善了本参考域的统计功效，但固定晚间PQ仍不能稳定给出唯一候选，说明评分改进与输入设计应结合。
90个补充完整AC归档含900个episode、1,296,000个快照；独立核对全部收敛且有限，KCL/功率平衡最大残差分别为$1.06\times10^{-13}$/$6.55\times10^{-13}$pu；竞争分数与校准真候选切片重算一致。
比较两法本身不授权在看到当前Y后挑选最愿给单例的方法；实际部署应事前固定方法，或把方法索引也纳入同时校准，才能保留同一5\%上界。
竞争差同时消去了所有模型都拟合很差的共同信息，因此不能替代模型适配检查。后续应将绝对适配性和相对区分作联合校准；本次未在追看结果后增加第三种选择规则。

\section{更现实的单向设备与波形优化}
补充研究固定为已成功预约的低压端口5、00--06窗口，保证域仅含以下五个预冻结动作。
每模板200个新校准episode、100个新测试episode；采用不同于主研究和初步四动作检查的新seed域。
现实ON--OFF要求驻留0或至少1.4kW、1min指令延迟、1kW/min爬坡；连续细调基线仍为理想接口。
被动$Q=0.2P$不是独立无功控制，只给一个组合响应方向。
表中未决含空集合和多候选集合，所有运行保留在分母。
\begin{center}\small
\begin{tabular}{L{5.7cm}rrrr}\toprule
同6h/9kWh方案 & 唯一正确\% & 覆盖\% & 错误\% & 未决\%\\\midrule
同能量恒功率 & 0.0 & 95.0 & 0.0 & 100.0 \\
连续细调ABBA & 82.3 & 94.0 & 0.0 & 17.7 \\
最小功率/延迟ON--OFF & 80.3 & 94.0 & 0.0 & 19.7 \\
ON--OFF且$Q=0.2P$ & 90.0 & 94.0 & 0.0 & 10.0 \\
稀疏码事前优化 & 87.7 & 94.0 & 0.0 & 12.3 \\
\bottomrule\end{tabular}
\end{center}
稀疏码在24个目标槽中安排6个高功率槽，使用0/6kW及25\%目标占空比。
先生成256个候选码，只以实际平均P在固定$[1,t]$干扰空间外的能量选优，随后才校准和测试。
有限码本可执行性、净能量与初末SOC逐条检查；这证明码本内事前改进，不是连续波形全局最优。
选中的高目标槽为第5--7及18--20槽，即00点开始时的01:00--01:45与04:15--05:00；其余目标为0。6kW为目标值，实际执行还包含1min延迟、爬坡和归零，辨识使用归档的真实15min平均值。
\begin{center}\footnotesize
\begin{tabular}{L{7.2cm}rrr}\toprule
配对比较 & 差值(pp) & 描述性95\%区间 & 本补充族Holm $p$\\\midrule
同能量恒功率 → 连续细调ABBA & +82.3 & [75.7, 88.3] & $1.62\times10^{-27}$ \\
连续细调ABBA → 最小功率/延迟ON--OFF & -2.0 & [-3.7, -0.7] & $0.03125$ \\
最小功率/延迟ON--OFF → ON--OFF且$Q=0.2P$ & +9.7 & [6.0, 13.7] & $1.43\times10^{-6}$ \\
最小功率/延迟ON--OFF → 稀疏码事前优化 & +7.3 & [4.0, 11.3] & $2.90\times10^{-4}$ \\
\bottomrule\end{tabular}
\end{center}
每模板的精确二项覆盖区间与完整100个seed簇记录保存于附加结果；三个模板不能当300个独立同分布样本。
本次主动方案经验覆盖为94\%，相应每模板100次的Clopper--Pearson区间为$[87.40\%,97.77\%]$；如实保留这一低于名义95\%的样本结果。
边际校准定理不要求一个已实现校准集的每个有限测试比例都达到95\%，这里也不能用零错误唯一来替代真模板覆盖。
新增波形重新校准，未沿用主60动作库的阈值。这给出了提升单向EV信息量的可实施方向，避免把更大的双向幅值误写成算法优势。

\section{失配与细步误差：保证何时停止适用}
\subsection{参考分布失准}
将全部实际线路R/X提高10\%，却仍使用原阻抗模板和原阈值，每尺度三模板各100次新测试。低压名义模板保留率92.33\%、中压97.00\%；这里真连续参数已不在模板中，所谓保留只指原结构标签，不能称为真模型95\%覆盖。

再固定已规划动作、评分与阈值，在30个全新seed簇上仅反转实际电压噪声的时段顺序；每簇三个模板，共90条记录。低压联合PQ在同seed参考条件下唯一正确/保留均96.67\%，反转后均6.67\%；中压联合PQ的保留率从93.33\%降至6.67\%。大量空集合表明原参考域不再适用；不能用原阈值继续声称95\%保证。
择时具有工程意义的前提是事前量测质量信息可信；部署应独立监测分布变化、扩大有效误差域或重新校准。用新增测试结果反向修改旧阈值，不能修复原试验的确认性。

\subsection{独立物理核对}
主研究保存150个完整训练/校准/测试归档，3780个episode共5,443,200个平衡非线性AC快照。全部收敛且电压有限；最终KCL最大残差$7.25\times10^{-13}$pu，欧姆定律残差$5.55\times10^{-17}$pu，复功率平衡残差$2.32\times10^{-12}$pu，电压范围$[0.98493,1.00091]$pu。
这是对保存的完整电流、电压、功率及损耗的重新计算；没有仅凭求解器收敛标志宣布物理正确。数百万相关快照仍只提供每尺度200个独立测试seed簇。


\subsection{AC平均的误差}
24个功率值来自连续轨迹精确平均，而非理想阶跃。非线性意味着
\[
 f_{\rm AC}(\overline P,\overline Q)\ne
 \overline{f_{\rm AC}(P(t),Q(t))},\qquad
 \overline{V^2}-(\overline V)^2=\operatorname{Var}_{\rm slot}(V).
\]
在两种五末端物理树、端口3的8个运行上固定槽内背景，分别用1min与0.5min中点积分。
下表给0.5min结果的平方电压降最大差及步长变化；单位为平方电压pu。
\begin{center}\small
\begin{tabular}{lrr}\toprule
控制 & 平均输入AC与平均响应之差 & 1min到0.5min积分变化\\\midrule
同能量恒功率 & $2.04\times10^{-8}$ & $1.05\times10^{-9}$ \\
单向ABBA P & $1.46\times10^{-7}$ & $3.58\times10^{-9}$ \\
单向PQ & $2.86\times10^{-6}$ & $1.21\times10^{-8}$ \\
双向PQ & $5.13\times10^{-6}$ & $2.17\times10^{-8}$ \\
\bottomrule\end{tabular}
\end{center}
该有限算例误差检查不是全允许域的严格上界；主大样本校准仍属于平均输入AC模拟law。
现场应明确仪表报告平均幅值、平均平方值还是其它RMS统计，再把执行/采样误差纳入统一偏差预算或直接以细步模型校准。

\section{WZZT衔接：局部准确不等于全网参数准确}
为避免固定阻抗模板掩盖参数辨识难度，另对两种五末端树各3个配对seed，使用被动24点加新24点。
两个预声明支持族均重新优化全部R/X，并枚举每条压缩边上的源前/后阻抗；选择只读量测，真值仅在最终评价中读取。
共30组拟合均完成两个支持族、全部源边LP，全部60批归档AC量测核对收敛且有限。
源幅值使用首端--分表估计及事前Q/P=0.35，本机制试验源幅值允许设0，故\textbf{不能把拟合误差解释为可靠置信认证}。

\begin{center}\small
\begin{tabular}{lrrrr}\toprule
新6h数据 & 端口3的R列误差\% & 端口3的X列误差\% & 全R误差\% & 全X误差\%\\\midrule
额外6h无车 & 19.2 & 87.8 & 22.0 & 102.2 \\
同能量恒功率 & 1.8 & 25.7 & 9.9 & 47.5 \\
单向ABBA P & 1.3 & 28.4 & 10.0 & 48.1 \\
单向PQ & 1.5 & 0.9 & 11.7 & 60.7 \\
双向PQ & 0.9 & 0.9 & 12.0 & 61.0 \\
\bottomrule\end{tabular}
\end{center}
所有族选择均正确，但这一小机制案例不足以给显著性或概率保证。
PQ明显补足被探测列的X方向；其余支路缺少独立输入，完整X矩阵仍可很差，且不同拟合解受源代理失配和病态条件影响。
因此下一地点应针对尚未认证支撑或未激励参数方向；全树结构完成后若还要求R/X精度，必须继续满足全参数设计的秩与条件数要求。

既有v3对同批AC数据优化contrast，R/X响应带平均全宽缩小4.17\%/7.66\%；这是固定曲线下的权重收益。
该凸半宽目标可写SOCP；当前数值实现采用可行解检查与FWL回退，未把SLSQP状态当全局最优证据。
前轮新增的独立数值工具给出精确存储浮点点积的泄漏界：$10^{-10}$量级干扰残差乘$10^9$量级系数可以造成约0.1偏差，原始容差本身不能证明覆盖。

\section{优化下界、LP对偶及合法拒绝}
\subsection{为什么必须用全局下界排除替代结构}
对结构$h$的完整允许参数域$\Theta_h$定义共同源原模型损失
$J_h=\inf_{\theta\in\Theta_h}L(Y;h,\theta)$，令$J_h^{\rm box}$为同域的源误差盒松弛。
\begin{proposition}[损失下界到错误拒绝控制]\label{prop:rejection}
若在概率至少$1-\alpha$的事件上真实$(h_\star,\theta_\star)$满足$\theta_\star\in\Theta_{h_\star}$且$L(Y;h_\star,\theta_\star)\le\tau$，并有经过验证的全域下界$\underline J_h\le J_h^{\rm box}$，则用$\underline J_h>\tau$拒绝$h$时，错误拒绝$h_\star$的概率至多$\alpha$。
\end{proposition}
\begin{proof}
在该事件上，命题\ref{prop:sourcebox}给
$\underline J_{h_\star}\le J_{h_\star}^{\rm box}\le J_{h_\star}\le L(Y;h_\star,\theta_\star)\le\tau$，所以拒绝事件只能出现在其补集。
\end{proof}
阈值条件需要对真参数所在域有效；求到一个低损失可行解只给$J_h$上界，不能满足上述排除方向。
若位置或其它离散分支域为$\mathcal E_h$，则
\[
 J_h^{\rm box}=\min_{e\in\mathcal E_h}J_{he}^{\rm box},\qquad
 \underline J_h=\min_{e\in\mathcal E_h}\underline J_{he}\le J_h^{\rm box}.
\]
未完成/失败分支须保留为未知（例如下界$-\infty$），不能静默删除；只有有效不可行性证明才可将其视为$+\infty$。
一个固定支持族内的所有源边枚举仍未覆盖其它支持族。

\subsection{LP乘子的严格下界意义}
\begin{proposition}[含对偶驻点残差的拉格朗日下界]\label{prop:lpbound}
对$\min c^\top x$、$Ax\ge b$、$\ell\le x\le u$，任取$\lambda\ge0$，置$r=c-A^\top\lambda$。若盒的支撑下界有限，则
\begin{equation}\label{eq:lpbound}
 D(\lambda)=b^\top\lambda+\sum_j\min\{r_j\ell_j,r_ju_j\}
 \le v_{\rm LP}\le v_{\rm MIP},
\end{equation}
右侧MIP为对同一个LP添加整数限制的模型。含无穷端点时用$\inf_{\ell\le x\le u}r^\top x$解释，不将$0\cdot\infty$作数值乘法。
\end{proposition}
\begin{proof}
任何原LP可行$x$都满足
$c^\top x=\lambda^\top Ax+r^\top x\ge\lambda^\top b+\inf_{[\ell,u]}r^\top x$。
对$x$取下确界得到第一不等式；整数可行域是LP可行域的子集，给第二个不等式。
\end{proof}
不需要将小残差当成0，也不需要下界乘子最优。存储LP的精确有理点积及最终向下舍入可保护该标量界，但不能保护未包含的物理参数/量化误差。

若$A,c,\ell,u$固定且$b(y)=b_0+By$，\emph{产生同一个$D$的同一组乘子}给
\begin{align}\label{eq:rhsbound}
 v_{\rm LP}(y+\Delta y)&\ge D(\lambda;y)+(B^\top\lambda)^\top\Delta y,\\
 \inf_{|\Delta y_i|\le\epsilon_i}v_{\rm LP}(y+\Delta y)
 &\ge D(\lambda;y)-\sum_i|(B^\top\lambda)_i|\epsilon_i.
\end{align}
这是代入$b(y+\Delta y)$后再取盒下确界的直接推论。最优对偶达到$v_{\rm LP}$时，该仿射下界为支持超平面；价格可能不唯一，不能无条件声称可微。
PQ或“源允许幅度乘未知$g$”改变$A$时，须重新计算残差校正或重解模型，不能只沿用RHS价格。

\begin{center}\small
\begin{tabular}{L{3.3cm}L{5.6cm}L{5.7cm}}\toprule
对象 & 严格用途 & 范围限制\\\midrule
完整整数松弛LP & 同模型MIP目标的全局下界 & 分数路径不等于合法拓扑\\
固定全部整数LP & 固定分支的最优值/敏感性 & 分支下界不能直接下界所有分支之最小值\\
分支定界搜索 & 合法全局bound及incumbent gap & bound不是对偶向量或正确性概率\\
\bottomrule\end{tabular}
\end{center}
例如$\min\{z:z\in\mathbb Z_+,z\ge b\}$在$b>0$的值为$\lceil b\rceil$，LP松弛值为$b$；前者是阶梯函数，不能把LP价格1解释为整数值函数处处导数。
固定整数后的价格只属于该分支\cite{gurobi,scipy}。

\textbf{当前接口限制：}\path{lp_dual.py}对乘子作符号投影，必要时统一缩小，再校正对偶驻点残差，已实现严格标量下界。
但导出的\path{target_rhs_marginals}仍是原始求解器边际，并非必然与校正下界匹配；严格应用式\eqref{eq:rhsbound}还需导出实际生成$D$的修正乘子。
历史有限差分检查只能支持其测试点的局部数值敏感性，不能填补这个接口缺口。

\section{更强置信目标：已推导但尚未实现}
\subsection{完整复合域的e-process}
令$\mathcal G_{r-1}$为历史，行动$a_r$及预测密度$q_r$在看到$Y_r$前确定，所有密度相对同一支配测度$\nu$。
\begin{proposition}[可选停止的错误拒绝控制]\label{prop:eprocess}
假设$q_r\ge0$且$\int q_r\,d\nu=1$，存在可测点态上界
\[
 \overline p_{h,r}(y)\ge\sup_{\theta\in\Theta_h}
 p_\theta(y\mid\mathcal G_{r-1},a_r).
\]
分母在真密度为正处须为正；全null零概率集合上的比值可任意定义。则对每个$\theta\in\Theta_h$，
\[
 \mathcal E_{h,0}=1,\qquad
 \mathcal E_{h,r}=\prod_{s=1}^r\frac{q_s(Y_s)}{\overline p_{h,s}(Y_s)}
\]
是非负超鞅，且$\Pr_\theta(\sup_{r\ge1}\mathcal E_{h,r}\ge1/\alpha)\le\alpha$。
\end{proposition}
\begin{proof}
给定历史，倍率的期望为
$\int q_r(y)p_\theta(y\mid\mathcal G_{r-1},a_r)/\overline p_{h,r}(y)\,d\nu(y)\le\int q_r\,d\nu=1$。
乘以前一时刻非负资本得到超鞅性质；应用Ville不等式。该不等式对每个允许真$\theta$都成立。
\end{proof}
这与universal inference及时间一致推断相连\cite{universal,howard}。
共同已知$d$维Gaussian协方差$\Sigma_r$下，若
\[
 \underline L_{h,r}(y)\le\inf_{\theta\in\Theta_h}
 \tfrac12\|y-\mu_\theta(a_r)\|_{\Sigma_r^{-1}}^2,
\]
则$\overline p_{h,r}(y)=(2\pi)^{-d/2}|\Sigma_r|^{-1/2}\exp[-\underline L_{h,r}(y)]$为合法分母，因为指数函数随损失递减。
局部AC可行损失是上界，方向相反；未知协方差还需处理归一化常数。
逐批对$\theta$取上确界较保守，但正是此逐步证明支持的构造。L1-LP界不能不经分布/尺度推导直接当Gaussian平方损失界；本路线尚无实现或功效试验。

\subsection{选择性风险：输出之后的群体错误率}
给定已经完成的训练和共形校准，冻结$M$套逐episode预测/弃权规则。
用\emph{另一个独立同分布、带真标签的风险校准集}，令规则$a$输出数为$n_a$、错误数为$k_a$，其部署law下选择性风险为$\rho_a=\Pr(\text{错误}\mid\text{输出})$，要求输出概率为正。
\begin{proposition}[独立风险校准后的安全规则选择]\label{prop:selective}
取单侧Clopper--Pearson上界
\[
 U_a=\operatorname{Beta}^{-1}(1-\delta/M;k_a+1,n_a-k_a),\quad0\le k_a<n_a;
\]
$n_a=0$或$k_a=n_a$时置$U_a=1$。则
$\Pr_{\rm riskcal}(\forall a:\rho_a\le U_a)\ge1-\delta$。
因此查看风险校准结果后选择任意$U_a\le\alpha$的规则，可获得相应law下$\rho_a\le\alpha$的$1-\delta$校准保证。
\end{proposition}
\begin{proof}
对每个固定逐episode规则，独立采样使给定$n_a$时$k_a\sim\operatorname{Bin}(n_a,\rho_a)$；二项尾反演给单侧覆盖至少$1-\delta/M$。
对$n_a$积分，再对规则并集即可；规则之间可以相关。
\end{proof}
当$k_a=0$，$U_a=1-(\delta/M)^{1/n_a}$，所以还需$n_a\ge\lceil\log(\delta/M)/\log(1-\alpha)\rceil$个\emph{输出样本}。
不能用该风险集重拟合规则，也不能对整个样本批按名额联动挑选输出后仍直接套上述二项条件。
这与Learn then Test的风险检验路线相接\cite{ltt}。若部署拓扑混合权重未知，可分别对每个$h$控制$\rho_{ah}$，以$\delta/(M|\mathcal H|)$分配；在校准好事件上，任意混合的选择性风险为
\[
 \rho_a^{\rm mix}=\frac{\sum_h\omega_h\Pr_h(\text{输出})\rho_{ah}}
 {\sum_h\omega_h\Pr_h(\text{输出})}\le\alpha,
\]
前提是分母为正，且所有有正输出概率的类均已验证$\rho_{ah}\le\alpha$。代价是每类需要足够独立输出样本。
当前三模板共享seed的300条记录不能作300个独立二项样本；这一选择性路线未实现，也仍不提供单个具体天气/节点输出的后验概率。

\section{共同源与历史联合拟合：三末端的严格机制验证}
\label{sec:localv6}
\textbf{本节为2026年10月11日新增的自行推导与独立小算例。}
它检验一个具体问题：源幅值允许逐时变化时，保留共同空间响应和历史共享线路参数，能否保留同电量EV调制的辨识收益？
本节采用精确线性潮流工作模型、三个预声明候选和平方电压降误差；它不是前述四案例AC外包络的替换，也不是完整WZZT搜索、源位置未知或现场性能验证。

\subsection{五项理论及本节使用范围}
\begin{center}\small
\begin{tabular}{L{3.0cm}L{11.6cm}}\toprule
理论 & 直觉及具体作用\\\midrule
树响应核 & 共享根路径决定电压响应。一个局部分叉可以由一列响应区分，不要求全部PQ方向满秩\cite{cavraro}。\\
共同源的集合几何 & 同一源在同一时刻只产生一个共同幅值；逐输出盒允许了源无法产生的空间变化。\\
剖面拟合与共享历史 & 各候选重新拟合允许参数，但线路参数不能随数据批次重置；源幅值仍可逐时变化。\\
模型判别设计 & 让竞争解释充分重拟合后再评价实验；T型设计是选动作的依据，不自动给正确概率\cite{dette}。本节固定两条曲线验证机制，尚未实现选点择时闭环。\\
集合辨识与条件排除 & 只有对整个声明候选参数域有效的正下界才能排除；真实误差在允许域内时保留真候选\cite{cerone2014}。\\\bottomrule
\end{tabular}
\end{center}
本节使用$L_\infty$额外相容性松弛，便于给出精确分离常数；不是声称现有WZZT的$L_1$目标已经改为此目标。

\subsection{物理结构、变量与精确LP}
根0先经公共干线到节点$v$，未计量源接在$v$；在$v$下，三个观测末端中的两端再共享一条分叉支路。
候选$h\in\{12,13,23\}$分别表示成组的两端。令$z_h$为该两端的三维指示向量，
\begin{equation}\label{eq:v6kernel}
 R_h(\theta)=a\mathbf1\mathbf1^\top+bz_hz_h^\top+\operatorname{diag}(c_1,c_2,c_3),
 \qquad\theta=(a,b,c_1,c_2,c_3).
\end{equation}
$a$是公共干线权重，$b$是分叉支路权重，$c_i$是末端支路权重，均已吸收平方电压降系数2。
功率以10 kW为基准，$y_i=V_0^2-V_i^2$以pu$^2$表示。令$Q=0$，总末端有功$p_t$已包含EV，
\begin{equation}\label{eq:v6model}
 y_t=R_h(\theta)p_t+a(\bar s_t+\delta s_t)\mathbf1+e_t,
 \quad |\delta s_t|\le B_t,\quad \|e_t\|_\infty\le\epsilon.
\end{equation}
源位置固定于公共干线末端是本节额外的明确前提；$\delta s_t$不要求平滑、独立、零均值，也不要求与充电曲线不相关。
其名义幅值$\bar s_t$外给。所有候选的允许参数相同：
\[
 a\in[.01,.05],\quad b\in[.005,.03],\quad c_i\in[.01,.12],
 \qquad \bar s_t=.12,\quad B_t=.10.
\]
因此全部允许实际源$\bar s_t+\delta s_t$非负，源误差可正可负。

\begin{proposition}[本局部源模型的精确线性化]\label{prop:v6exactlp}
固定候选$h$，以$w_t=a\delta s_t$代换，则$|w_t|\le B_ta$与原源约束精确等价。
候选的最小额外相容性松弛为下列LP的最优值：
\begin{equation}\label{eq:v6lp}
 \tau_h^\star=\min_{\theta,w,\tau\ge0}\tau\quad\text{s.t.}\quad
 \begin{cases}
 |y_{it}-[R_h(\theta)p_t]_i-a\bar s_t-w_t|\le\epsilon+\tau,\quad\forall i,t,\\
 -B_ta\le w_t\le B_ta,\quad\theta\text{在上述参数盒内}.
 \end{cases}
\end{equation}
\end{proposition}
\begin{proof}
原可行点给出$w_t=a\delta s_t$；反向由$a\ge.01>0$取$\delta s_t=w_t/a$即得原幅值约束。
式\eqref{eq:v6kernel}对$\theta$线性，$p_t,\bar s_t,B_t$都是给定输入，因此拆开绝对值后全为线性约束。
该等价利用了源响应只有$a\mathbf1$且幅值逐时自由；一般未知源位置、多个不同路径权重与同一幅值相乘时，不能据此宣称仍为精确LP。
\end{proof}

\paragraph{两组公平消融。}
源模型比较：共同源使用同一$w_t$，误差盒改用各节点独立$w_{it}$且$|w_{it}|\le B_ta$；其它输入和参数域完全一致。
历史比较：joint在旧、新两批共享$\theta$；reset在同一份旧、新观测上给两批各自的$\theta$，共用一个$\tau$，所以其目标恰为两批单独最小松弛的最大值。
new-only仅使用新批，是另一种信息量较少的对照，不与reset混称。

\begin{proposition}[包含关系与真实候选保留]\label{prop:v6nest}
在相同观测、误差预算和参数域下，
\[
 \tau_{h,\mathrm{shared}}^\star\ge\tau_{h,\mathrm{box}}^\star,
 \qquad \tau_{h,\mathrm{joint}}^\star\ge\tau_{h,\mathrm{reset}}^\star.
\]
若真实候选在三候选域内，真实参数和源满足上述域，且$\|e\|_\infty\le\epsilon$，则真候选在共同源joint中有$\tau^\star=0$。
\end{proposition}
\begin{proof}
共同源可行点可取$w_{it}=w_t$嵌入误差盒；joint可行点可复制参数嵌入reset。较大可行域的最小目标不会变大。
将真实参数、真实$w_t$及$\tau=0$代入即可证明真候选保留。两个不等式都不保证严格增益。
\end{proof}
这里是给定误差事件上的确定性结论；本节不额外给出95\%覆盖或拓扑后验概率。若将来误差事件有经过验证的概率覆盖，该覆盖才能传递至保留候选及其公共结构。

\subsection{非平凡精确分离：完整证明与可行见证}
生成树为$h=12$，真实$\theta=(.02,.01,.06,.06,.06)$。末端背景为$(.2,.2,.2)$，EV接在末端1，令
\[
 p_t=(.2+u_t,.2,.2),\qquad |\delta s_t|\le.02.
\]
这个较小的实际源变化幅度仅用于构造和检验算例；推断仍允许完整的$B=.10$。
旧批12点取$u=0$。新批24个15 min槽：恒充取$u=.15$；调制取$u=0,.3,.3,0$重复6次。
两者分别对应1.5 kW恒充和0/3 kW单向充电，均6 h、9 kWh网侧电量。

\begin{proposition}[两种错误分叉的精确无噪声距离]\label{prop:v6margin}
设全部用于联合拟合的$u_t\in[0,L]$，两端点均出现，$L\in\{.15,.30\}$。
在式\eqref{eq:v6lp}取无噪声数据、$\epsilon=0$时，即使允许竞争者重新拟合全部$a,b,c_i$与逐时共同源，也有
\begin{equation}\label{eq:v6exactmargin}
 d_{13}:=\tau_{13}^\star=\frac{.015L}{4},\qquad
 d_{23}:=\tau_{23}^\star=\frac{.010L}{4}.
\end{equation}
\end{proposition}
\begin{proof}
取空间对比$y_2-y_3$，公共干线和源均抵消。真树的该对比对$u$的斜率为$.01$；候选13的斜率为$-b'$，候选23的斜率为0。
若所有坐标残差不超过$d$，两个时点、两个坐标组成的差分误差最多为$4d$。
故$(.01+b')L\le4d$或$.01L\le4d$。利用$b'\ge.005$得到式\eqref{eq:v6exactmargin}的下界。

下面给出达到下界的显式参数，而非以求解器状态代替存在证明：
\begin{center}\small
\begin{tabular}{lrrrrrl}\toprule
候选 & $a'$ & $b'$ & $c'_1$ & $c'_2$ & $c'_3$ & $w'_t$\\\midrule
13 & .0225 & .005 & .0625 & $.0725+.01875L$ & $.0425-.01875L$ & $-.0003+.02\delta s_t$\\
23 & .0250 & .005 & .0650 & $.0550+.01250L$ & $.0350-.01250L$ & $-.0006+.02\delta s_t$\\\bottomrule
\end{tabular}
\end{center}
两行参数均在声明盒内。由$|\delta s_t|\le.02$，分别有$|w'_t|\le.0007<.00225=a'B$与$|w'_t|\le.0010<.0025=a'B$。
直接代入式\eqref{eq:v6model}，真响应减去候选响应分别为
\[
 r^{13}_t=(0,\ .0075(u_t-L/2),\ -.0075(u_t-L/2)),
 \qquad
 r^{23}_t=(0,\ .005(u_t-L/2),\ -.005(u_t-L/2)).
\]
端点处最大绝对残差恰为$.015L/4$和$.010L/4$，与下界一致。证明覆盖全部允许竞争参数，不只是固定参数模板。
\end{proof}

\paragraph{历史与调制分别增加了什么。}
joint恒充使用$u=0,.15$，所以$d_{13}=.0005625,d_{23}=.000375$；joint调制使用$u=0,.3$，两者分别为$.001125,.000750$。
因此三候选内最小错误结构距离由$.000375$增至$.000750$，在这个精确模型中加倍。
对恒充的new-only或reset，每一批只有一个$u$值，可各自改选末端阻抗精确拟合，距离为0；调制新批自身已有两个水平，故本算例再加旧批不提高其无噪声距离。
\textbf{后一项零增益同样是结果：历史有用不意味着每次追加历史都严格增强当前实验。}

\begin{proposition}[有界噪声及误差盒的零排除]\label{prop:v6noisebox}
记$d_h$为上述无噪声最小距离。若真实观测噪声不超过$\rho$而LP允许误差$\epsilon$，则
\[
 [d_h-\rho-\epsilon]_+\le\tau_h^\star\le[d_h+\rho-\epsilon]_+.
\]
特别地，$\rho=\epsilon$时$[d_h-2\epsilon]_+\le\tau_h^\star\le d_h$。
而在本算例$L\le.30$、$\rho\le\epsilon$时，逐输出误差盒能使13、23两错误候选均有$\tau^\star=0$。
\end{proposition}
\begin{proof}
到同一预测集合的$L_\infty$距离随数据改变至多$\rho$；允许$\epsilon$后，额外松弛为距离减$\epsilon$的正部，得到第一组界。
对误差盒，将上个证明的残差各坐标分别加入$w'_t$即可精确匹配无噪声数据。
候选13所需$|w'_{it}|$最多为$.0007+.001125=.001825<.00225$；候选23最多为$.0010+.000750=.001750<.0025$。
实际量测噪声由$\epsilon$吸收，故两者仍以$\tau=0$可行。源盒可以吞掉共同源不能产生的一正一负结构残差。
\end{proof}
取$\epsilon=10^{-5}$，恒充和调制的三候选最小分离下界分别为$.000355$与$.000730$。
由于调制下界还大于恒充上界$.000375$，收益在整个声明噪声盒内严格成立，不依赖一个有利随机种子。
本证明中的噪声是平方电压降噪声；不能据此宣称原始电压幅值高斯噪声经过平方后仍满足相同分布。

\paragraph{正结果不只发生在一个真参数点。}
固定上述输入和任意$|\delta s_t|\le.02$，若真实五个参数各自变化不超过$\eta$，无噪声响应的最大坐标变化不超过
\[
 (1.04+.70+.50)\eta=2.24\eta.
\]
三项分别界住公共干线承载的总功率（含实际源）、分叉下游功率和单末端功率。
到固定候选预测集合的距离以及两错误候选距离的最小值均为1-Lipschitz函数。
故有界噪声下，调制的最小错误候选松弛至少为$.000750-2.24\eta-2\epsilon$，恒充至多为$.000375+2.24\eta$。
当$4.48\eta+2\epsilon<.000375$时，严格收益仍成立。
例如$\eta=5\times10^{-5},\epsilon=10^{-5}$时，全部真参数仍位于允许盒内部，调制下界$.000618$大于恒充上界$.000487$。
这是一个非空真参数邻域的结论；没有声称所有允许真参数上都有相同收益，也不由连续性自动推出误差盒仍精确零拟合。

\subsection{对照测试、负结果与未完成的推广}
本节数值实现、冻结配置、原始观测、逐次求解界与独立复核保存在\path{research_experiments/mobile_ev_probing/results_v6/}。
实现入口为\path{local_shared_source_v6.py}。全部204次LP包括192次主/无噪声核验和12次负例求解，不是204个独立现场实验。
正式核验使用事先冻结的seed 1201、1202、1203；每个seed的两条曲线使用相同源轨迹、相同槽位噪声和相同旧数据。
实际源误差均匀生成于$[-.02,.02]$，平方压降噪声生成于$[-.8\epsilon,.8\epsilon]$，其中$\epsilon=10^{-5}$；推断仅接收观测和声明预算，不接收真拓扑、真阻抗或真实源幅值。
种子71的无噪声核验先与解析式对照，浮点观测舍入由显式微小误差预算覆盖；正式三个种子不用于修改动作或参数域。
探索两策略的舍入预算分别约为$3.382\times10^{-18}$与$3.412\times10^{-18}$，并非逐位相同；正式评估的误差预算严格相同，独立归档审计将两者分开处理。
下表列出解析无噪声距离及一致的数值排除模式，距离单位为$10^{-4}$ pu$^2$。

\begin{center}\small
\begin{tabular}{llrrrr}\toprule
源模型 & 数据/参数处理 & 恒充$d_{13}$ & 恒充$d_{23}$ & 调制$d_{13}$ & 调制$d_{23}$\\\midrule
共同源 & joint：旧新共享 & 5.625 & 3.750 & 11.250 & 7.500\\
共同源 & reset：同数据分批参数 & 0 & 0 & 11.250 & 7.500\\
共同源 & new-only：仅新批 & 0 & 0 & 11.250 & 7.500\\
共同源 & old-only：仅旧批 & 0 & 0 & 0 & 0\\
逐输出盒 & 上述四种处理 & 0 & 0 & 0 & 0\\\bottomrule
\end{tabular}
\end{center}
这里的零表示存在允许解释或本数值审计未获得正下界，不是将微小负数截成收益。
误差盒零值及共同源恒充分批零值另有显式可行构造；数值零不单独承担不可辨识性证明。
对恒充单批$u=U$，13可取$a'=a,b'=b,c'=(.06,.06+.01(2+U/.2),.06-.01(2+U/.2))$；
23可取$c'_1=.06+.01(.4+U)/(.2+U),c'_2=.06+.01U/.2,c'_3=.04$，均在$U=0,.15$时精确拟合。

有噪声的joint共同源最优额外松弛在三个种子中的范围如下。
区间端点按所有已验证下界的最小值与上界的最大值向外显示，\textbf{是跨种子的求解值范围，不是统计置信区间}。
\begin{center}\small
\begin{tabular}{lrr}\toprule
错误候选 & 恒充$\tau^\star$范围（$10^{-4}$） & 调制$\tau^\star$范围（$10^{-4}$）\\\midrule
13 & [5.5839, 5.5865] & [11.1860, 11.2118]\\
23 & [3.7089, 3.7115] & [7.4360, 7.4618]\\\bottomrule
\end{tabular}
\end{center}
三个种子中，共同源joint的恒充和调制都排除两个错误候选；reset/new-only恒充均不排除，调制均排除两个；源盒均不排除。
因此在本噪声水平下，调制的收益表现为\emph{分离余量提高}，不是把已经100\%的这组三种子排除率再提高。
三个种子只作数值复核，不承担跨网络显著性或现场可靠率结论；式\eqref{eq:v6exactmargin}及命题\ref{prop:v6noisebox}承担本局部模型的严格结论。

\paragraph{数值证书与设备审计。}
每个LP保存HiGHS状态、原始可行性残差、KKT诊断、系数外向舍入预算、精确有理数校正的外LP对偶下界，以及在原始有理数系数模型上修复和核验的可行上界。
下界来自外LP，仍下界原问题；上界必须验证原模型，不能把仅对放宽LP可行的点当作原模型上界。
204次均完成，最大上下界间隙小于$2.40\times10^{-16}$，最大原始浮点行违约小于$4.86\times10^{-17}$，没有求解失败，真候选均保留。
另有6项独立测试通过，含8次小LP、解析达界见证、共同源与EV完全相关时的有界噪声检验及退化情形；与204次研究求解分列，不合并为独立样本数。
数值有理数审计针对归档的二进制输入，解析十进制公式与其微小表示误差分开记录。
程序仅为本协议设$\tau\le1$；全部已验证原模型可行上界小于$.001126$，故该上限不改变本次最优值，不能把此固定上限当作任意新数据的通用接口保证。

EV采用固定接入点、单向3 kW额定功率、无独立Q；50 kWh电池从SOC $.4$开始，充电效率固定$.95$，允许SOC $[.2,.9]$。
两策略网侧电量均9 kWh、电池增能均8.55 kWh、最终SOC均$.571$。
本次只检查15 min槽平均功率、槽间变化和SOC，不模拟启动延迟、切换瞬态、连续时间爬坡或AC电压动态；由$\sqrt{1-y}$计算的电压范围也只属于线性模型范围检查。

\paragraph{保留的负结果与适用边界。}
\begin{enumerate}
\item \textbf{误差过大仍未决：}保持同一真实分叉、全部历史和调制，令允许压降误差$\epsilon=.002$、实际观测无附加随机噪声，三候选在共同源和误差盒中全部保留（6次LP）。更宽误差域覆盖真实数据，但无法提供有用排除。
\item \textbf{分叉坍缩仍不可辨：}令真实$b=0$，同时允许域改为$b\in[0,.03]$，其它条件不变；三候选都可表示同一星形响应，全部保留（6次LP）。正分叉权重是本分离证明的必要条件之一。
独立测试还取严格正的$b=10^{-6}$、同步降低候选下界而保持$\epsilon=10^{-5}$，错误13仍未决：正分叉本身不够，结构间隔还必须大于量测误差。
\item \textbf{新增历史可能零收益：}调制新批已经覆盖两个极值，本例joint与new-only无噪声距离相同；不能把样本数增加解释成必然的信息增益。
\item \textbf{源位置推广尚未证明：}本例源恰为共同模态，$y_2-y_3$能够完全消源。源在争议支路、位置未知或多源时不再自动成立，需对实际允许位置重新拟合和证明；四案例未知源域全部未决的记录继续保留。
\item \textbf{未检验选点择时：}端口1、旧新时窗和两条曲线预先固定，不能据此宣称已有自动地点/时段选择优势。当前$[P;Q]$仍全方向秩亏，但有限局部候选可分；这不意味着全网R/X已可恢复。
\end{enumerate}
本次已得到的结果是：\emph{存在连续阻抗域、非零时变源误差和有界噪声下可严格验证的局部EV收益，逐输出误差盒确实可能完全丢失该收益。}
下一步应先把源位置从公共干线扩展到少量有工程依据的位置，检查这个正间隔是否保留；若间隔消失，应报告需要额外源位置或量测信息，不能只更换设计指标。

\section{可拓展方向、实施顺序与复现}
\begin{center}\small
\begin{tabular}{L{2.2cm}L{6.1cm}L{6.3cm}}\toprule
优先级 & 下一步 & 可检验成果\\\midrule
1 & 用公开实车轨迹及目标车验收标定Q(P)、启动、跟踪和SOC效率 & 形成设备动作域；控制误差传播进置信半宽\\
1 & 将单向有限稀疏码、时地T分数与WZZT未决支撑连接 & 同服务与作业成本减少未决；对手重拟合全部R/X/源路径\\
2 & 三相、配变、线路电流/电压及接入限制 & 事前可行域筛选和执行后全过程审计，而非仅设备SOC通过\\
2 & 含未知阻抗与源位置的双侧最坏分离 & 内层全局下界或可核验松弛，报告零分离及遗漏域\\
2 & 绝对适配性与竞争差联合校准；独立校准选择性风险 & 明确误排、输出条件错误率及检验族，报告有效样本不足\\
3 & 多车/多轮时窗和路径MILP，联合SOC状态 & 正确累计旅行、等待、能量和信息相关性；重新定义策略episode\\
3 & 全局损失下界支撑复合e-process & 明确条件似然、拒绝风险及无法停止的情况\\
\bottomrule\end{tabular}
\end{center}
若随机充电码在背景发生前生成且背景不响应EV，可用成对Rademacher码使干扰相关项平均抵消；源反馈、控制内生性和设备跟踪破坏该排除条件。
随机化/工具变量是值得研究的第四种设计思想，尚未完成现场因果验证。
脉冲压缩文献提供编码和相关处理背景\cite{pcp}，其周期级动态等效模型不能直接用来证明本项目6h静态拓扑恢复。

研究代码位于\path{research_experiments/mobile_ev_probing/}，新增设备、候选置信集合、时地研究、源拟合衔接及数值泄漏模块。
2026年10月6日数值版本记录\textbf{429项测试通过}。
10月8日复核数学与四案例历史证据，修复候选置信接口的语义元数据；当日该模块\textbf{16项针对性测试通过}（包含3项新增语义检查），未重跑数值研究，也不将16项重复计入历史总数。
10月9日起新增的观测LP、同时误差盒、独立节点AC对抗测试、四案例协议及报告证据门控检查共\textbf{108项测试通过}；与历史计数分列，不合并宣传为一个新测试总数。
本轮四例新实验、正式与探索性结果区分、逐候选上下界、原始量测、代码指纹和启动前设计修订归档于\path{results_v5/}，详见本节前的同域验证及\path{README_v5.md}。
原始AC输出、实际P/Q、带噪电压、逐次选择、置信集合、失败和代码指纹均归档于\path{results_v4/}；既有v2/v3数据保留。
本文件是用户已打开\path{docs/mobile_ev_probing_report_v2.tex}的原位续修，文件名中的v2不表示内容仍为旧版。

在仓库根目录使用Topo环境Python，设置\path{PYTHONPATH=.;rnj_wzzt_core}，按顺序运行模块：
\begin{quote}\small\ttfamily
research\_experiments.mobile\_ev\_probing.action\_study\_v4\\
research\_experiments.mobile\_ev\_probing.capability\_study\_v4\\
research\_experiments.mobile\_ev\_probing.sampling\_audit\_v4\\
research\_experiments.mobile\_ev\_probing.wzzt\_bridge\_v4\\
research\_experiments.mobile\_ev\_probing.confidence\_compare\_v4\\
research\_experiments.mobile\_ev\_probing.confidence\_score\_study\_v4\\
research\_experiments.mobile\_ev\_probing.verify\_v4\\
research\_experiments.mobile\_ev\_probing.update\_open\_report\_v4
\end{quote}
主研究还需运行结果目录内保存的时段失配与AC独立审计脚本，具体命令见研究README。
报告生成器要求所有正式结果完整，只原位更新当前源文件；用户要求的PDF由同一源文件另行编译，不替换编辑器标签。
已有证据支持声明模型内的有效机制及条件结论；新LP覆盖已枚举树的完整R/X盒，尚未覆盖任意未知拓扑。
现场分布、设备实测、网络安全和缩减树之外的物理细节仍须新增证据。

\clearpage
\begin{thebibliography}{99}\small
\bibitem{cavraro} Cavraro G., Kekatos V. Graph Algorithms for Topology Identification using Power Grid Probing. IEEE Control Systems Letters, 2018. \url{https://arxiv.org/abs/1803.04506}.
\bibitem{taheri} Taheri S. et al. An MILP Approach for Distribution Grid Topology Identification Using Inverter Probing. \url{https://arxiv.org/abs/2004.02370}.
\bibitem{bhela} Bhela S., Kekatos V., Veeramachaneni S. Smart Inverter Grid Probing for Learning Loads: Part II---Probing Injection Design. \url{https://arxiv.org/abs/1806.08836}.
\bibitem{hotz} Hotz M. et al. Topology Estimation in Low Voltage Grids Using Wallbox Charging Data Recordings. CIRED, 2022. \url{https://doi.org/10.1049/icp.2022.0823}.
\bibitem{doe} U.S. DOE FEMP. Bidirectional Charging and Electric Vehicles for Mobile Storage. \url{https://www.energy.gov/cmei/femp/bidirectional-charging-and-electric-vehicles-mobile-storage}.
\bibitem{evse} U.S. DOE FEMP. Electric Vehicle Supply Equipment Infrastructure: Federal Fleet Training. \url{https://www.energy.gov/cmei/femp/articles/electric-vehicle-supply-equipment-infrastructure-federal-fleet-training}.
\bibitem{saxena} Saxena S. et al. Field Testing of Residential Bidirectional Electric Vehicle Charger for Power System Applications. 2023. \url{https://arxiv.org/abs/2308.13624}.
\bibitem{esser} Esser M. et al. High-temporal-resolution dataset of uni-, bidirectional, and dynamic electric vehicle charging profiles. Scientific Data, 2025, 12:1192. \url{https://doi.org/10.1038/s41597-025-05524-5}. 数据：\url{https://zenodo.org/records/14065331}.
\bibitem{dette} Dette H., Titoff S. Optimal discrimination designs. Annals of Statistics, 2009, 37:2056--2082. \url{https://arxiv.org/abs/0908.1912}.
\bibitem{armstrong} Armstrong T., Kolesár M. Optimal Inference in a Class of Regression Models. \url{https://arxiv.org/abs/1511.06028}.
\bibitem{wilks} Wilks S. Determination of Sample Sizes for Setting Tolerance Limits. Annals of Mathematical Statistics, 1941. \url{https://doi.org/10.1214/aoms/1177731788}.
\bibitem{vovk} Vovk V. Conditional validity of inductive conformal predictors. PMLR, 2012, 25:475--490. \url{https://proceedings.mlr.press/v25/vovk12.html}.
\bibitem{universal} Wasserman L., Ramdas A., Balakrishnan S. Universal inference. PNAS, 2020, 117:16880--16890. \url{https://arxiv.org/abs/1912.11436}.
\bibitem{howard} Howard S. et al. Time-uniform, nonparametric, nonasymptotic confidence sequences. Annals of Statistics, 2021. \url{https://arxiv.org/abs/1810.08240}.
\bibitem{pcp} Piaquadio N. et al. Pulse Compression Probing for Tracking Distribution Feeder Models. 2023. \url{https://arxiv.org/abs/2305.19465}.
\bibitem{gurobi} Gurobi. How do I retrieve the dual (Pi) values for a MIP problem? \url{https://support.gurobi.com/hc/en-us/articles/360034305272-How-do-I-retrieve-the-dual-Pi-values-for-a-MIP-problem}.
\bibitem{scipy} SciPy. linprog(method='highs'): marginals and solver diagnostics. \url{https://docs.scipy.org/doc/scipy/reference/optimize.linprog-highs.html}.
\bibitem{margin} Aleksandrova M., Chertov R. Impact of model-agnostic nonconformity functions on efficiency of conformal classifiers: an extensive study. PMLR, 2021, 152:151--170. \url{https://proceedings.mlr.press/v152/aleksandrova21a.html}.
\bibitem{ltt} Angelopoulos A. N. et al. Learn then Test: Calibrating Predictive Algorithms to Achieve Risk Control. \url{https://arxiv.org/abs/2110.01052}.
\bibitem{cerone2014} Cerone V., Lasserre J.-B., Piga D., Regruto D. A Unified Framework for Solving a General Class of Conditional and Robust Set-Membership Estimation Problems. IEEE Transactions on Automatic Control, 2014, 59(11):2897--2909. \url{https://doi.org/10.1109/TAC.2014.2351695}. 已读作者稿：\url{https://arxiv.org/abs/1408.0532v1}.
\bibitem{scott2014} Scott J. K., Findeisen R., Braatz R. D., Raimondo D. M. Input design for guaranteed fault diagnosis using zonotopes. Automatica, 2014, 50:1580--1589. \url{https://doi.org/10.1016/j.automatica.2014.03.016}.
\bibitem{ding2019} Ding Y., Harirchi F., Yong S. Z., Jacobsen E., Ozay N. Optimal Input Design for Affine Model Discrimination with Applications in Intention-Aware Vehicles. 本次采用2019年1月2日v3作者稿：\url{https://arxiv.org/abs/1702.01112v3}.
\bibitem{farivar} Farivar M., Low S. H. Branch Flow Model: Relaxations and Convexification---Part I. IEEE Transactions on Power Systems, 2013. 本次核对v4作者稿：\url{https://arxiv.org/abs/1204.4865v4}.
\bibitem{evminimum} ABB E-mobility. smart charge@home: Frequently asked questions. 已核对最低6 A及单/三相配置说明：\url{https://e-mobility.abb.com/en/smart-chargehome-frequently-asked-questions}. Schneider Electric. eMobility Infrastructure Commissioning Guide, Setting Energy Management. 其最小电流设定表列出单相及三相6 A：\url{https://productinfo.se.com/emobility-infrastructure-commissioning/evsolcg001-emobility-infrastructure-commissioning-guide/English/BM_EVSOLCG001EN_eMobility_Infrastructure_Commissioning_Guide_DD01016549.xml/$/EVSOLCG001EN_EV_SetEnergyManagement_DD01037747}. 以上是制造商技术文档核查，不是对任意车型的统一能力保证。
\end{thebibliography}
\end{document}
